%% inexe.tex
%% A display-math-style linguistic example environment.
%% Load with %% inexe.tex
%% A display-math-style linguistic example environment.
%% Load with %% inexe.tex
%% A display-math-style linguistic example environment.
%% Load with %% inexe.tex
%% A display-math-style linguistic example environment.
%% Load with \input{inexe} in your preamble.
%%
%% Configuration (set these BEFORE \input{inexe}):
%%
%%   \def\inexedelimiter{parens}   % or: brackets, period
%%   \newlength{\inexelabelsep}
%%   \setlength{\inexelabelsep}{1em}
%%
%% Counter sharing with gb4e:
%%   If gb4e is loaded, inexe automatically shares gb4e's \exx counter.
%%
%% Usage:
%%   \begin{inexe}
%%     \ex{Consider sub-examples:
%%       \begin{inxlist}     % a. b. c.  (default, like xlist)
%%         \xi First sub-example.
%%         \xij[*] Bad sub-example.
%%         \xisn Unnumbered sub-example.
%%       \end{inxlist}}
%%     \ex{Roman sub-examples:
%%       \begin{inxlisti}    % i. ii. iii.  (like xlisti)
%%         \xi First.
%%         \xi Second.
%%       \end{inxlisti}}
%%   \end{inexe}
%%
%% Inxlist variants (all support \xi, \xij, \xisn, \xisnj):
%%   inxlist   a. b. c.      (mirrors xlist/xlista)
%%   inxlisti  i. ii. iii.   (mirrors xlisti)
%%   inxlistn  1. 2. 3.      (mirrors xlistn)
%%   inxlistA  A. B. C.      (mirrors xlistA)
%%   inxlistI  I. II. III.   (mirrors xlistI)

%% ---------------------------------------------------------------
%% Guard against double-loading
%% ---------------------------------------------------------------
\expandafter\ifx\csname inexe@loaded\endcsname\relax
\expandafter\def\csname inexe@loaded\endcsname{}
\else
  \expandafter\endinput
\fi

\makeatletter

%% ---------------------------------------------------------------
%% User-configurable defaults
%% (Only set if not already defined by the user before \input.)
%% ---------------------------------------------------------------

\@ifundefined{inexedelimiter}{\def\inexedelimiter{parens}}{}

\@ifundefined{inexelabelsep}{%
  \newlength{\inexelabelsep}%
  \setlength{\inexelabelsep}{1em}%
}{}

%% ---------------------------------------------------------------
%% Label formatting
%% ---------------------------------------------------------------

\def\inexe@fmt#1{%
  \ifx\inexedelimiter\@inexe@parens (#1)\else
  \ifx\inexedelimiter\@inexe@brackets [#1]\else
  \ifx\inexedelimiter\@inexe@period #1.\else
  (#1)%
  \fi\fi\fi
}
\def\@inexe@parens{parens}
\def\@inexe@brackets{brackets}
\def\@inexe@period{period}

%% ---------------------------------------------------------------
%% Counter
%% ---------------------------------------------------------------

\@definecounter{inexecnt}

\def\inexe@stepcounter{%
  \refstepcounter{inexecnt}%
  \protected@edef\@currentlabel{\arabic{inexecnt}}%
}
\def\inexe@thecounter{\arabic{inexecnt}}

\AtBeginDocument{%
  \@ifundefined{c@exx}{}{%
    \def\inexe@stepcounter{%
      \refstepcounter{exx}%
      \setcounter{inexecnt}{\value{exx}}%
      \protected@edef\@currentlabel{\arabic{exx}}%
    }%
    \def\inexe@thecounter{\arabic{exx}}%
    \def\theinexecnt{\arabic{exx}}%
  }%
}

%% ---------------------------------------------------------------
%% Core display line
%% ---------------------------------------------------------------

\newlength{\inexe@savedleft}
\newlength{\inexe@labelwidth}
\newlength{\inexe@itemsep}
\setlength{\inexe@itemsep}{3\p@ plus2\p@}

\long\def\inexe@line#1#2{%
  \if@inexe@firstline
    \@inexe@firstlinefalse
  \else
    \par
  \fi
  \noindent
  \hspace*{-\inexe@savedleft}%
  \noindent
  \hbox to \dimexpr\linewidth+\inexe@savedleft\relax{%
    \hbox to \inexe@labelwidth{%
      \hfil
      \ifx\relax#1\relax
      \else
        \inexe@fmt{#1}%
      \fi
    }%
    \hspace{\inexelabelsep}%
    \parbox[t]{\dimexpr
        \linewidth + \inexe@savedleft
        - \inexe@labelwidth
        - \inexelabelsep
      \relax}{#2}%
    \hss
  }%
  \par
  \vspace{\inexe@itemsep}%
}

%% ---------------------------------------------------------------
%% Item commands
%% ---------------------------------------------------------------

\long\def\inexe@ex#1{%
  \inexe@stepcounter
  \inexe@line{\inexe@thecounter}{#1}%
}

\long\def\inexe@exj[#1]#2{%
  \inexe@stepcounter
  \inexe@line{\inexe@thecounter}{#1\hspace{.5ex}#2}%
}

\long\def\inexe@sn#1{%
  \inexe@line{\relax}{#1}%
}

\long\def\inexe@snj[#1]#2{%
  \inexe@line{\relax}{#1\hspace{.5ex}#2}%
}

\long\def\inexe@exi#1#2{%
  \inexe@line{#1}{#2}%
}

%% ---------------------------------------------------------------
%% inxlist — inline sub-example list for use inside \ex{...}
%%
%% Lives inside the \parbox of \ex{...}. Uses \halign so \xi takes
%% unbraced content (everything to the next \xi/\xij/\xisn/\xisnj
%% or \end{inxlist}), matching gb4e's xlist feel.
%% Labels are always a. b. c. style (gb4e xlist default).
%%
%% Usage:
%%   \ex{Consider the following:
%%     \begin{inxlist}
%%       \xi First sub-example.
%%       \xi Second sub-example.
%%       \xij[*] Bad sub-example.
%%       \xisn Unnumbered sub-example.
%%     \end{inxlist}}
%%
%% Commands inside inxlist:
%%   \xi          lettered sub-example (unbraced)
%%   \xij[jdg]    with judgement mark
%%   \xisn        unnumbered sub-example
%%   \xisnj[jdg]  unnumbered with judgement
%% ---------------------------------------------------------------

\@definecounter{inxcnt}
% \p@inxcnt: prefix for \ref — produces "3b" style cross-references.
% Defined as the current inexe label (arabic) so sub-example labels
% inherit the parent example number.
\def\p@inxcnt{\theinexecnt}
% \theinexecnt: the current top-level example number, used as prefix.
% Points to exx if gb4e loaded, else inexecnt. Defined at \begin{document}
% (after gb4e has loaded) via \AtBeginDocument below — placeholder for now.
\def\theinexecnt{\arabic{inexecnt}}
% Save the standard \@item before any package (e.g. paralist) redefines it.
\let\@orignalitem\@item

\newlength{\inxlist@labelwidth}
\newlength{\inxlist@labelsep}
\setlength{\inxlist@labelsep}{.75em}
\newlength{\inxlist@itemsep}
\setlength{\inxlist@itemsep}{1.5\p@ plus\p@}

% \theinxcnt: formats the sub-example counter value (letter/roman/etc.)
% Set locally by \inxlist@open to match the current list style.
\def\theinxcnt{\alph{inxcnt}}
%% Shared setup for all inxlist variants.
%% #1 = counter format command, e.g. \alph, \roman, \arabic, \Alph, \Roman
%% #2 = sample string for label width, e.g. "m." or "iii." or "99."
\def\inxlist@open#1#2{%
  \setcounter{inxcnt}{0}%
  \def\inxlist@counterfmt{#1}%
  % Set \theinxcnt so \refstepcounter{inxcnt} records the right label.
  \def\theinxcnt{#1{inxcnt}}%
  \settowidth{\inxlist@labelwidth}{#2}%
  \edef\inxlist@savedlistdepth{\the\@listdepth}%
  \edef\inxlist@saveditemdepth{\the\@itemdepth}%
  \let\@item\@orignalitem
  \begin{list}{}{%
    \labelwidth=\inxlist@labelwidth
    \labelsep=\inxlist@labelsep
    \leftmargin=\inxlist@labelwidth
    \advance\leftmargin\inxlist@labelsep
    \itemindent\z@
    \listparindent\z@
    \topsep=\inxlist@itemsep
    \itemsep=\inxlist@itemsep
    \parsep\z@
  }%
  \def\xi{%
    \refstepcounter{inxcnt}%
    \item[\inxlist@counterfmt{inxcnt}.]%
  }%
  \def\xij[##1]{%
    \refstepcounter{inxcnt}%
    \item[\inxlist@counterfmt{inxcnt}.]##1\hspace{.5ex}%
  }%
  \def\xisn{\item[]}%
  \def\xisnj[##1]{\item[]##1\hspace{.5ex}}%
}

\def\inxlist@close{%
  \end{list}%
  \@listdepth=\inxlist@savedlistdepth\relax
  \@itemdepth=\inxlist@saveditemdepth\relax
  \ignorespacesafterend
}

% inxlist  — alphabetic:  a. b. c.   (mirrors gb4e's xlist / xlista)
\newenvironment{inxlist}
  {\inxlist@open{\alph}{m.}}
  {\inxlist@close}

% inxlisti — roman:       i. ii. iii.  (mirrors gb4e's xlisti)
\newenvironment{inxlisti}
  {\inxlist@open{\roman}{iii.}}
  {\inxlist@close}

% inxlistn — arabic:      1. 2. 3.   (mirrors gb4e's xlistn)
\newenvironment{inxlistn}
  {\inxlist@open{\arabic}{9.}}
  {\inxlist@close}

% inxlistA — Alphabetic:  A. B. C.   (mirrors gb4e's xlistA)
\newenvironment{inxlistA}
  {\inxlist@open{\Alph}{M.}}
  {\inxlist@close}

% inxlistI — Roman:       I. II. III.  (mirrors gb4e's xlistI)
\newenvironment{inxlistI}
  {\inxlist@open{\Roman}{III.}}
  {\inxlist@close}

%% ---------------------------------------------------------------
%% inexe environment
%% ---------------------------------------------------------------

\newif\if@inexe@firstline

\newenvironment{inexe}{%
  \inexe@savedleft=\@totalleftmargin
  \advance\inexe@savedleft\leftskip
  \ifnum\value{inexecnt}<9
    \settowidth{\inexe@labelwidth}{\inexe@fmt{9}}%
  \else\ifnum\value{inexecnt}<99
    \settowidth{\inexe@labelwidth}{\inexe@fmt{99}}%
  \else
    \settowidth{\inexe@labelwidth}{\inexe@fmt{999}}%
  \fi\fi
  \@inexe@firstlinetrue
  \leavevmode\newline\noindent
  \let\ex\inexe@ex
  \let\exj\inexe@exj
  \let\sn\inexe@sn
  \let\snj\inexe@snj
  \let\exi\inexe@exi
}{%
  \vspace{-\inexe@itemsep}%
  \vspace{\belowdisplayskip}%
  \ignorespacesafterend
}

\makeatother

\endinput
 in your preamble.
%%
%% Configuration (set these BEFORE %% inexe.tex
%% A display-math-style linguistic example environment.
%% Load with \input{inexe} in your preamble.
%%
%% Configuration (set these BEFORE \input{inexe}):
%%
%%   \def\inexedelimiter{parens}   % or: brackets, period
%%   \newlength{\inexelabelsep}
%%   \setlength{\inexelabelsep}{1em}
%%
%% Counter sharing with gb4e:
%%   If gb4e is loaded, inexe automatically shares gb4e's \exx counter.
%%
%% Usage:
%%   \begin{inexe}
%%     \ex{Consider sub-examples:
%%       \begin{inxlist}     % a. b. c.  (default, like xlist)
%%         \xi First sub-example.
%%         \xij[*] Bad sub-example.
%%         \xisn Unnumbered sub-example.
%%       \end{inxlist}}
%%     \ex{Roman sub-examples:
%%       \begin{inxlisti}    % i. ii. iii.  (like xlisti)
%%         \xi First.
%%         \xi Second.
%%       \end{inxlisti}}
%%   \end{inexe}
%%
%% Inxlist variants (all support \xi, \xij, \xisn, \xisnj):
%%   inxlist   a. b. c.      (mirrors xlist/xlista)
%%   inxlisti  i. ii. iii.   (mirrors xlisti)
%%   inxlistn  1. 2. 3.      (mirrors xlistn)
%%   inxlistA  A. B. C.      (mirrors xlistA)
%%   inxlistI  I. II. III.   (mirrors xlistI)

%% ---------------------------------------------------------------
%% Guard against double-loading
%% ---------------------------------------------------------------
\expandafter\ifx\csname inexe@loaded\endcsname\relax
\expandafter\def\csname inexe@loaded\endcsname{}
\else
  \expandafter\endinput
\fi

\makeatletter

%% ---------------------------------------------------------------
%% User-configurable defaults
%% (Only set if not already defined by the user before \input.)
%% ---------------------------------------------------------------

\@ifundefined{inexedelimiter}{\def\inexedelimiter{parens}}{}

\@ifundefined{inexelabelsep}{%
  \newlength{\inexelabelsep}%
  \setlength{\inexelabelsep}{1em}%
}{}

%% ---------------------------------------------------------------
%% Label formatting
%% ---------------------------------------------------------------

\def\inexe@fmt#1{%
  \ifx\inexedelimiter\@inexe@parens (#1)\else
  \ifx\inexedelimiter\@inexe@brackets [#1]\else
  \ifx\inexedelimiter\@inexe@period #1.\else
  (#1)%
  \fi\fi\fi
}
\def\@inexe@parens{parens}
\def\@inexe@brackets{brackets}
\def\@inexe@period{period}

%% ---------------------------------------------------------------
%% Counter
%% ---------------------------------------------------------------

\@definecounter{inexecnt}

\def\inexe@stepcounter{%
  \refstepcounter{inexecnt}%
  \protected@edef\@currentlabel{\arabic{inexecnt}}%
}
\def\inexe@thecounter{\arabic{inexecnt}}

\AtBeginDocument{%
  \@ifundefined{c@exx}{}{%
    \def\inexe@stepcounter{%
      \refstepcounter{exx}%
      \setcounter{inexecnt}{\value{exx}}%
      \protected@edef\@currentlabel{\arabic{exx}}%
    }%
    \def\inexe@thecounter{\arabic{exx}}%
    \def\theinexecnt{\arabic{exx}}%
  }%
}

%% ---------------------------------------------------------------
%% Core display line
%% ---------------------------------------------------------------

\newlength{\inexe@savedleft}
\newlength{\inexe@labelwidth}
\newlength{\inexe@itemsep}
\setlength{\inexe@itemsep}{3\p@ plus2\p@}

\long\def\inexe@line#1#2{%
  \if@inexe@firstline
    \@inexe@firstlinefalse
  \else
    \par
  \fi
  \noindent
  \hspace*{-\inexe@savedleft}%
  \noindent
  \hbox to \dimexpr\linewidth+\inexe@savedleft\relax{%
    \hbox to \inexe@labelwidth{%
      \hfil
      \ifx\relax#1\relax
      \else
        \inexe@fmt{#1}%
      \fi
    }%
    \hspace{\inexelabelsep}%
    \parbox[t]{\dimexpr
        \linewidth + \inexe@savedleft
        - \inexe@labelwidth
        - \inexelabelsep
      \relax}{#2}%
    \hss
  }%
  \par
  \vspace{\inexe@itemsep}%
}

%% ---------------------------------------------------------------
%% Item commands
%% ---------------------------------------------------------------

\long\def\inexe@ex#1{%
  \inexe@stepcounter
  \inexe@line{\inexe@thecounter}{#1}%
}

\long\def\inexe@exj[#1]#2{%
  \inexe@stepcounter
  \inexe@line{\inexe@thecounter}{#1\hspace{.5ex}#2}%
}

\long\def\inexe@sn#1{%
  \inexe@line{\relax}{#1}%
}

\long\def\inexe@snj[#1]#2{%
  \inexe@line{\relax}{#1\hspace{.5ex}#2}%
}

\long\def\inexe@exi#1#2{%
  \inexe@line{#1}{#2}%
}

%% ---------------------------------------------------------------
%% inxlist — inline sub-example list for use inside \ex{...}
%%
%% Lives inside the \parbox of \ex{...}. Uses \halign so \xi takes
%% unbraced content (everything to the next \xi/\xij/\xisn/\xisnj
%% or \end{inxlist}), matching gb4e's xlist feel.
%% Labels are always a. b. c. style (gb4e xlist default).
%%
%% Usage:
%%   \ex{Consider the following:
%%     \begin{inxlist}
%%       \xi First sub-example.
%%       \xi Second sub-example.
%%       \xij[*] Bad sub-example.
%%       \xisn Unnumbered sub-example.
%%     \end{inxlist}}
%%
%% Commands inside inxlist:
%%   \xi          lettered sub-example (unbraced)
%%   \xij[jdg]    with judgement mark
%%   \xisn        unnumbered sub-example
%%   \xisnj[jdg]  unnumbered with judgement
%% ---------------------------------------------------------------

\@definecounter{inxcnt}
% \p@inxcnt: prefix for \ref — produces "3b" style cross-references.
% Defined as the current inexe label (arabic) so sub-example labels
% inherit the parent example number.
\def\p@inxcnt{\theinexecnt}
% \theinexecnt: the current top-level example number, used as prefix.
% Points to exx if gb4e loaded, else inexecnt. Defined at \begin{document}
% (after gb4e has loaded) via \AtBeginDocument below — placeholder for now.
\def\theinexecnt{\arabic{inexecnt}}
% Save the standard \@item before any package (e.g. paralist) redefines it.
\let\@orignalitem\@item

\newlength{\inxlist@labelwidth}
\newlength{\inxlist@labelsep}
\setlength{\inxlist@labelsep}{.75em}
\newlength{\inxlist@itemsep}
\setlength{\inxlist@itemsep}{1.5\p@ plus\p@}

% \theinxcnt: formats the sub-example counter value (letter/roman/etc.)
% Set locally by \inxlist@open to match the current list style.
\def\theinxcnt{\alph{inxcnt}}
%% Shared setup for all inxlist variants.
%% #1 = counter format command, e.g. \alph, \roman, \arabic, \Alph, \Roman
%% #2 = sample string for label width, e.g. "m." or "iii." or "99."
\def\inxlist@open#1#2{%
  \setcounter{inxcnt}{0}%
  \def\inxlist@counterfmt{#1}%
  % Set \theinxcnt so \refstepcounter{inxcnt} records the right label.
  \def\theinxcnt{#1{inxcnt}}%
  \settowidth{\inxlist@labelwidth}{#2}%
  \edef\inxlist@savedlistdepth{\the\@listdepth}%
  \edef\inxlist@saveditemdepth{\the\@itemdepth}%
  \let\@item\@orignalitem
  \begin{list}{}{%
    \labelwidth=\inxlist@labelwidth
    \labelsep=\inxlist@labelsep
    \leftmargin=\inxlist@labelwidth
    \advance\leftmargin\inxlist@labelsep
    \itemindent\z@
    \listparindent\z@
    \topsep=\inxlist@itemsep
    \itemsep=\inxlist@itemsep
    \parsep\z@
  }%
  \def\xi{%
    \refstepcounter{inxcnt}%
    \item[\inxlist@counterfmt{inxcnt}.]%
  }%
  \def\xij[##1]{%
    \refstepcounter{inxcnt}%
    \item[\inxlist@counterfmt{inxcnt}.]##1\hspace{.5ex}%
  }%
  \def\xisn{\item[]}%
  \def\xisnj[##1]{\item[]##1\hspace{.5ex}}%
}

\def\inxlist@close{%
  \end{list}%
  \@listdepth=\inxlist@savedlistdepth\relax
  \@itemdepth=\inxlist@saveditemdepth\relax
  \ignorespacesafterend
}

% inxlist  — alphabetic:  a. b. c.   (mirrors gb4e's xlist / xlista)
\newenvironment{inxlist}
  {\inxlist@open{\alph}{m.}}
  {\inxlist@close}

% inxlisti — roman:       i. ii. iii.  (mirrors gb4e's xlisti)
\newenvironment{inxlisti}
  {\inxlist@open{\roman}{iii.}}
  {\inxlist@close}

% inxlistn — arabic:      1. 2. 3.   (mirrors gb4e's xlistn)
\newenvironment{inxlistn}
  {\inxlist@open{\arabic}{9.}}
  {\inxlist@close}

% inxlistA — Alphabetic:  A. B. C.   (mirrors gb4e's xlistA)
\newenvironment{inxlistA}
  {\inxlist@open{\Alph}{M.}}
  {\inxlist@close}

% inxlistI — Roman:       I. II. III.  (mirrors gb4e's xlistI)
\newenvironment{inxlistI}
  {\inxlist@open{\Roman}{III.}}
  {\inxlist@close}

%% ---------------------------------------------------------------
%% inexe environment
%% ---------------------------------------------------------------

\newif\if@inexe@firstline

\newenvironment{inexe}{%
  \inexe@savedleft=\@totalleftmargin
  \advance\inexe@savedleft\leftskip
  \ifnum\value{inexecnt}<9
    \settowidth{\inexe@labelwidth}{\inexe@fmt{9}}%
  \else\ifnum\value{inexecnt}<99
    \settowidth{\inexe@labelwidth}{\inexe@fmt{99}}%
  \else
    \settowidth{\inexe@labelwidth}{\inexe@fmt{999}}%
  \fi\fi
  \@inexe@firstlinetrue
  \leavevmode\newline\noindent
  \let\ex\inexe@ex
  \let\exj\inexe@exj
  \let\sn\inexe@sn
  \let\snj\inexe@snj
  \let\exi\inexe@exi
}{%
  \vspace{-\inexe@itemsep}%
  \vspace{\belowdisplayskip}%
  \ignorespacesafterend
}

\makeatother

\endinput
):
%%
%%   \def\inexedelimiter{parens}   % or: brackets, period
%%   \newlength{\inexelabelsep}
%%   \setlength{\inexelabelsep}{1em}
%%
%% Counter sharing with gb4e:
%%   If gb4e is loaded, inexe automatically shares gb4e's \exx counter.
%%
%% Usage:
%%   \begin{inexe}
%%     \ex{Consider sub-examples:
%%       \begin{inxlist}     % a. b. c.  (default, like xlist)
%%         \xi First sub-example.
%%         \xij[*] Bad sub-example.
%%         \xisn Unnumbered sub-example.
%%       \end{inxlist}}
%%     \ex{Roman sub-examples:
%%       \begin{inxlisti}    % i. ii. iii.  (like xlisti)
%%         \xi First.
%%         \xi Second.
%%       \end{inxlisti}}
%%   \end{inexe}
%%
%% Inxlist variants (all support \xi, \xij, \xisn, \xisnj):
%%   inxlist   a. b. c.      (mirrors xlist/xlista)
%%   inxlisti  i. ii. iii.   (mirrors xlisti)
%%   inxlistn  1. 2. 3.      (mirrors xlistn)
%%   inxlistA  A. B. C.      (mirrors xlistA)
%%   inxlistI  I. II. III.   (mirrors xlistI)

%% ---------------------------------------------------------------
%% Guard against double-loading
%% ---------------------------------------------------------------
\expandafter\ifx\csname inexe@loaded\endcsname\relax
\expandafter\def\csname inexe@loaded\endcsname{}
\else
  \expandafter\endinput
\fi

\makeatletter

%% ---------------------------------------------------------------
%% User-configurable defaults
%% (Only set if not already defined by the user before \input.)
%% ---------------------------------------------------------------

\@ifundefined{inexedelimiter}{\def\inexedelimiter{parens}}{}

\@ifundefined{inexelabelsep}{%
  \newlength{\inexelabelsep}%
  \setlength{\inexelabelsep}{1em}%
}{}

%% ---------------------------------------------------------------
%% Label formatting
%% ---------------------------------------------------------------

\def\inexe@fmt#1{%
  \ifx\inexedelimiter\@inexe@parens (#1)\else
  \ifx\inexedelimiter\@inexe@brackets [#1]\else
  \ifx\inexedelimiter\@inexe@period #1.\else
  (#1)%
  \fi\fi\fi
}
\def\@inexe@parens{parens}
\def\@inexe@brackets{brackets}
\def\@inexe@period{period}

%% ---------------------------------------------------------------
%% Counter
%% ---------------------------------------------------------------

\@definecounter{inexecnt}

\def\inexe@stepcounter{%
  \refstepcounter{inexecnt}%
  \protected@edef\@currentlabel{\arabic{inexecnt}}%
}
\def\inexe@thecounter{\arabic{inexecnt}}

\AtBeginDocument{%
  \@ifundefined{c@exx}{}{%
    \def\inexe@stepcounter{%
      \refstepcounter{exx}%
      \setcounter{inexecnt}{\value{exx}}%
      \protected@edef\@currentlabel{\arabic{exx}}%
    }%
    \def\inexe@thecounter{\arabic{exx}}%
    \def\theinexecnt{\arabic{exx}}%
  }%
}

%% ---------------------------------------------------------------
%% Core display line
%% ---------------------------------------------------------------

\newlength{\inexe@savedleft}
\newlength{\inexe@labelwidth}
\newlength{\inexe@itemsep}
\setlength{\inexe@itemsep}{3\p@ plus2\p@}

\long\def\inexe@line#1#2{%
  \if@inexe@firstline
    \@inexe@firstlinefalse
  \else
    \par
  \fi
  \noindent
  \hspace*{-\inexe@savedleft}%
  \noindent
  \hbox to \dimexpr\linewidth+\inexe@savedleft\relax{%
    \hbox to \inexe@labelwidth{%
      \hfil
      \ifx\relax#1\relax
      \else
        \inexe@fmt{#1}%
      \fi
    }%
    \hspace{\inexelabelsep}%
    \parbox[t]{\dimexpr
        \linewidth + \inexe@savedleft
        - \inexe@labelwidth
        - \inexelabelsep
      \relax}{#2}%
    \hss
  }%
  \par
  \vspace{\inexe@itemsep}%
}

%% ---------------------------------------------------------------
%% Item commands
%% ---------------------------------------------------------------

\long\def\inexe@ex#1{%
  \inexe@stepcounter
  \inexe@line{\inexe@thecounter}{#1}%
}

\long\def\inexe@exj[#1]#2{%
  \inexe@stepcounter
  \inexe@line{\inexe@thecounter}{#1\hspace{.5ex}#2}%
}

\long\def\inexe@sn#1{%
  \inexe@line{\relax}{#1}%
}

\long\def\inexe@snj[#1]#2{%
  \inexe@line{\relax}{#1\hspace{.5ex}#2}%
}

\long\def\inexe@exi#1#2{%
  \inexe@line{#1}{#2}%
}

%% ---------------------------------------------------------------
%% inxlist — inline sub-example list for use inside \ex{...}
%%
%% Lives inside the \parbox of \ex{...}. Uses \halign so \xi takes
%% unbraced content (everything to the next \xi/\xij/\xisn/\xisnj
%% or \end{inxlist}), matching gb4e's xlist feel.
%% Labels are always a. b. c. style (gb4e xlist default).
%%
%% Usage:
%%   \ex{Consider the following:
%%     \begin{inxlist}
%%       \xi First sub-example.
%%       \xi Second sub-example.
%%       \xij[*] Bad sub-example.
%%       \xisn Unnumbered sub-example.
%%     \end{inxlist}}
%%
%% Commands inside inxlist:
%%   \xi          lettered sub-example (unbraced)
%%   \xij[jdg]    with judgement mark
%%   \xisn        unnumbered sub-example
%%   \xisnj[jdg]  unnumbered with judgement
%% ---------------------------------------------------------------

\@definecounter{inxcnt}
% \p@inxcnt: prefix for \ref — produces "3b" style cross-references.
% Defined as the current inexe label (arabic) so sub-example labels
% inherit the parent example number.
\def\p@inxcnt{\theinexecnt}
% \theinexecnt: the current top-level example number, used as prefix.
% Points to exx if gb4e loaded, else inexecnt. Defined at \begin{document}
% (after gb4e has loaded) via \AtBeginDocument below — placeholder for now.
\def\theinexecnt{\arabic{inexecnt}}
% Save the standard \@item before any package (e.g. paralist) redefines it.
\let\@orignalitem\@item

\newlength{\inxlist@labelwidth}
\newlength{\inxlist@labelsep}
\setlength{\inxlist@labelsep}{.75em}
\newlength{\inxlist@itemsep}
\setlength{\inxlist@itemsep}{1.5\p@ plus\p@}

% \theinxcnt: formats the sub-example counter value (letter/roman/etc.)
% Set locally by \inxlist@open to match the current list style.
\def\theinxcnt{\alph{inxcnt}}
%% Shared setup for all inxlist variants.
%% #1 = counter format command, e.g. \alph, \roman, \arabic, \Alph, \Roman
%% #2 = sample string for label width, e.g. "m." or "iii." or "99."
\def\inxlist@open#1#2{%
  \setcounter{inxcnt}{0}%
  \def\inxlist@counterfmt{#1}%
  % Set \theinxcnt so \refstepcounter{inxcnt} records the right label.
  \def\theinxcnt{#1{inxcnt}}%
  \settowidth{\inxlist@labelwidth}{#2}%
  \edef\inxlist@savedlistdepth{\the\@listdepth}%
  \edef\inxlist@saveditemdepth{\the\@itemdepth}%
  \let\@item\@orignalitem
  \begin{list}{}{%
    \labelwidth=\inxlist@labelwidth
    \labelsep=\inxlist@labelsep
    \leftmargin=\inxlist@labelwidth
    \advance\leftmargin\inxlist@labelsep
    \itemindent\z@
    \listparindent\z@
    \topsep=\inxlist@itemsep
    \itemsep=\inxlist@itemsep
    \parsep\z@
  }%
  \def\xi{%
    \refstepcounter{inxcnt}%
    \item[\inxlist@counterfmt{inxcnt}.]%
  }%
  \def\xij[##1]{%
    \refstepcounter{inxcnt}%
    \item[\inxlist@counterfmt{inxcnt}.]##1\hspace{.5ex}%
  }%
  \def\xisn{\item[]}%
  \def\xisnj[##1]{\item[]##1\hspace{.5ex}}%
}

\def\inxlist@close{%
  \end{list}%
  \@listdepth=\inxlist@savedlistdepth\relax
  \@itemdepth=\inxlist@saveditemdepth\relax
  \ignorespacesafterend
}

% inxlist  — alphabetic:  a. b. c.   (mirrors gb4e's xlist / xlista)
\newenvironment{inxlist}
  {\inxlist@open{\alph}{m.}}
  {\inxlist@close}

% inxlisti — roman:       i. ii. iii.  (mirrors gb4e's xlisti)
\newenvironment{inxlisti}
  {\inxlist@open{\roman}{iii.}}
  {\inxlist@close}

% inxlistn — arabic:      1. 2. 3.   (mirrors gb4e's xlistn)
\newenvironment{inxlistn}
  {\inxlist@open{\arabic}{9.}}
  {\inxlist@close}

% inxlistA — Alphabetic:  A. B. C.   (mirrors gb4e's xlistA)
\newenvironment{inxlistA}
  {\inxlist@open{\Alph}{M.}}
  {\inxlist@close}

% inxlistI — Roman:       I. II. III.  (mirrors gb4e's xlistI)
\newenvironment{inxlistI}
  {\inxlist@open{\Roman}{III.}}
  {\inxlist@close}

%% ---------------------------------------------------------------
%% inexe environment
%% ---------------------------------------------------------------

\newif\if@inexe@firstline

\newenvironment{inexe}{%
  \inexe@savedleft=\@totalleftmargin
  \advance\inexe@savedleft\leftskip
  \ifnum\value{inexecnt}<9
    \settowidth{\inexe@labelwidth}{\inexe@fmt{9}}%
  \else\ifnum\value{inexecnt}<99
    \settowidth{\inexe@labelwidth}{\inexe@fmt{99}}%
  \else
    \settowidth{\inexe@labelwidth}{\inexe@fmt{999}}%
  \fi\fi
  \@inexe@firstlinetrue
  \leavevmode\newline\noindent
  \let\ex\inexe@ex
  \let\exj\inexe@exj
  \let\sn\inexe@sn
  \let\snj\inexe@snj
  \let\exi\inexe@exi
}{%
  \vspace{-\inexe@itemsep}%
  \vspace{\belowdisplayskip}%
  \ignorespacesafterend
}

\makeatother

\endinput
 in your preamble.
%%
%% Configuration (set these BEFORE %% inexe.tex
%% A display-math-style linguistic example environment.
%% Load with %% inexe.tex
%% A display-math-style linguistic example environment.
%% Load with \input{inexe} in your preamble.
%%
%% Configuration (set these BEFORE \input{inexe}):
%%
%%   \def\inexedelimiter{parens}   % or: brackets, period
%%   \newlength{\inexelabelsep}
%%   \setlength{\inexelabelsep}{1em}
%%
%% Counter sharing with gb4e:
%%   If gb4e is loaded, inexe automatically shares gb4e's \exx counter.
%%
%% Usage:
%%   \begin{inexe}
%%     \ex{Consider sub-examples:
%%       \begin{inxlist}     % a. b. c.  (default, like xlist)
%%         \xi First sub-example.
%%         \xij[*] Bad sub-example.
%%         \xisn Unnumbered sub-example.
%%       \end{inxlist}}
%%     \ex{Roman sub-examples:
%%       \begin{inxlisti}    % i. ii. iii.  (like xlisti)
%%         \xi First.
%%         \xi Second.
%%       \end{inxlisti}}
%%   \end{inexe}
%%
%% Inxlist variants (all support \xi, \xij, \xisn, \xisnj):
%%   inxlist   a. b. c.      (mirrors xlist/xlista)
%%   inxlisti  i. ii. iii.   (mirrors xlisti)
%%   inxlistn  1. 2. 3.      (mirrors xlistn)
%%   inxlistA  A. B. C.      (mirrors xlistA)
%%   inxlistI  I. II. III.   (mirrors xlistI)

%% ---------------------------------------------------------------
%% Guard against double-loading
%% ---------------------------------------------------------------
\expandafter\ifx\csname inexe@loaded\endcsname\relax
\expandafter\def\csname inexe@loaded\endcsname{}
\else
  \expandafter\endinput
\fi

\makeatletter

%% ---------------------------------------------------------------
%% User-configurable defaults
%% (Only set if not already defined by the user before \input.)
%% ---------------------------------------------------------------

\@ifundefined{inexedelimiter}{\def\inexedelimiter{parens}}{}

\@ifundefined{inexelabelsep}{%
  \newlength{\inexelabelsep}%
  \setlength{\inexelabelsep}{1em}%
}{}

%% ---------------------------------------------------------------
%% Label formatting
%% ---------------------------------------------------------------

\def\inexe@fmt#1{%
  \ifx\inexedelimiter\@inexe@parens (#1)\else
  \ifx\inexedelimiter\@inexe@brackets [#1]\else
  \ifx\inexedelimiter\@inexe@period #1.\else
  (#1)%
  \fi\fi\fi
}
\def\@inexe@parens{parens}
\def\@inexe@brackets{brackets}
\def\@inexe@period{period}

%% ---------------------------------------------------------------
%% Counter
%% ---------------------------------------------------------------

\@definecounter{inexecnt}

\def\inexe@stepcounter{%
  \refstepcounter{inexecnt}%
  \protected@edef\@currentlabel{\arabic{inexecnt}}%
}
\def\inexe@thecounter{\arabic{inexecnt}}

\AtBeginDocument{%
  \@ifundefined{c@exx}{}{%
    \def\inexe@stepcounter{%
      \refstepcounter{exx}%
      \setcounter{inexecnt}{\value{exx}}%
      \protected@edef\@currentlabel{\arabic{exx}}%
    }%
    \def\inexe@thecounter{\arabic{exx}}%
    \def\theinexecnt{\arabic{exx}}%
  }%
}

%% ---------------------------------------------------------------
%% Core display line
%% ---------------------------------------------------------------

\newlength{\inexe@savedleft}
\newlength{\inexe@labelwidth}
\newlength{\inexe@itemsep}
\setlength{\inexe@itemsep}{3\p@ plus2\p@}

\long\def\inexe@line#1#2{%
  \if@inexe@firstline
    \@inexe@firstlinefalse
  \else
    \par
  \fi
  \noindent
  \hspace*{-\inexe@savedleft}%
  \noindent
  \hbox to \dimexpr\linewidth+\inexe@savedleft\relax{%
    \hbox to \inexe@labelwidth{%
      \hfil
      \ifx\relax#1\relax
      \else
        \inexe@fmt{#1}%
      \fi
    }%
    \hspace{\inexelabelsep}%
    \parbox[t]{\dimexpr
        \linewidth + \inexe@savedleft
        - \inexe@labelwidth
        - \inexelabelsep
      \relax}{#2}%
    \hss
  }%
  \par
  \vspace{\inexe@itemsep}%
}

%% ---------------------------------------------------------------
%% Item commands
%% ---------------------------------------------------------------

\long\def\inexe@ex#1{%
  \inexe@stepcounter
  \inexe@line{\inexe@thecounter}{#1}%
}

\long\def\inexe@exj[#1]#2{%
  \inexe@stepcounter
  \inexe@line{\inexe@thecounter}{#1\hspace{.5ex}#2}%
}

\long\def\inexe@sn#1{%
  \inexe@line{\relax}{#1}%
}

\long\def\inexe@snj[#1]#2{%
  \inexe@line{\relax}{#1\hspace{.5ex}#2}%
}

\long\def\inexe@exi#1#2{%
  \inexe@line{#1}{#2}%
}

%% ---------------------------------------------------------------
%% inxlist — inline sub-example list for use inside \ex{...}
%%
%% Lives inside the \parbox of \ex{...}. Uses \halign so \xi takes
%% unbraced content (everything to the next \xi/\xij/\xisn/\xisnj
%% or \end{inxlist}), matching gb4e's xlist feel.
%% Labels are always a. b. c. style (gb4e xlist default).
%%
%% Usage:
%%   \ex{Consider the following:
%%     \begin{inxlist}
%%       \xi First sub-example.
%%       \xi Second sub-example.
%%       \xij[*] Bad sub-example.
%%       \xisn Unnumbered sub-example.
%%     \end{inxlist}}
%%
%% Commands inside inxlist:
%%   \xi          lettered sub-example (unbraced)
%%   \xij[jdg]    with judgement mark
%%   \xisn        unnumbered sub-example
%%   \xisnj[jdg]  unnumbered with judgement
%% ---------------------------------------------------------------

\@definecounter{inxcnt}
% \p@inxcnt: prefix for \ref — produces "3b" style cross-references.
% Defined as the current inexe label (arabic) so sub-example labels
% inherit the parent example number.
\def\p@inxcnt{\theinexecnt}
% \theinexecnt: the current top-level example number, used as prefix.
% Points to exx if gb4e loaded, else inexecnt. Defined at \begin{document}
% (after gb4e has loaded) via \AtBeginDocument below — placeholder for now.
\def\theinexecnt{\arabic{inexecnt}}
% Save the standard \@item before any package (e.g. paralist) redefines it.
\let\@orignalitem\@item

\newlength{\inxlist@labelwidth}
\newlength{\inxlist@labelsep}
\setlength{\inxlist@labelsep}{.75em}
\newlength{\inxlist@itemsep}
\setlength{\inxlist@itemsep}{1.5\p@ plus\p@}

% \theinxcnt: formats the sub-example counter value (letter/roman/etc.)
% Set locally by \inxlist@open to match the current list style.
\def\theinxcnt{\alph{inxcnt}}
%% Shared setup for all inxlist variants.
%% #1 = counter format command, e.g. \alph, \roman, \arabic, \Alph, \Roman
%% #2 = sample string for label width, e.g. "m." or "iii." or "99."
\def\inxlist@open#1#2{%
  \setcounter{inxcnt}{0}%
  \def\inxlist@counterfmt{#1}%
  % Set \theinxcnt so \refstepcounter{inxcnt} records the right label.
  \def\theinxcnt{#1{inxcnt}}%
  \settowidth{\inxlist@labelwidth}{#2}%
  \edef\inxlist@savedlistdepth{\the\@listdepth}%
  \edef\inxlist@saveditemdepth{\the\@itemdepth}%
  \let\@item\@orignalitem
  \begin{list}{}{%
    \labelwidth=\inxlist@labelwidth
    \labelsep=\inxlist@labelsep
    \leftmargin=\inxlist@labelwidth
    \advance\leftmargin\inxlist@labelsep
    \itemindent\z@
    \listparindent\z@
    \topsep=\inxlist@itemsep
    \itemsep=\inxlist@itemsep
    \parsep\z@
  }%
  \def\xi{%
    \refstepcounter{inxcnt}%
    \item[\inxlist@counterfmt{inxcnt}.]%
  }%
  \def\xij[##1]{%
    \refstepcounter{inxcnt}%
    \item[\inxlist@counterfmt{inxcnt}.]##1\hspace{.5ex}%
  }%
  \def\xisn{\item[]}%
  \def\xisnj[##1]{\item[]##1\hspace{.5ex}}%
}

\def\inxlist@close{%
  \end{list}%
  \@listdepth=\inxlist@savedlistdepth\relax
  \@itemdepth=\inxlist@saveditemdepth\relax
  \ignorespacesafterend
}

% inxlist  — alphabetic:  a. b. c.   (mirrors gb4e's xlist / xlista)
\newenvironment{inxlist}
  {\inxlist@open{\alph}{m.}}
  {\inxlist@close}

% inxlisti — roman:       i. ii. iii.  (mirrors gb4e's xlisti)
\newenvironment{inxlisti}
  {\inxlist@open{\roman}{iii.}}
  {\inxlist@close}

% inxlistn — arabic:      1. 2. 3.   (mirrors gb4e's xlistn)
\newenvironment{inxlistn}
  {\inxlist@open{\arabic}{9.}}
  {\inxlist@close}

% inxlistA — Alphabetic:  A. B. C.   (mirrors gb4e's xlistA)
\newenvironment{inxlistA}
  {\inxlist@open{\Alph}{M.}}
  {\inxlist@close}

% inxlistI — Roman:       I. II. III.  (mirrors gb4e's xlistI)
\newenvironment{inxlistI}
  {\inxlist@open{\Roman}{III.}}
  {\inxlist@close}

%% ---------------------------------------------------------------
%% inexe environment
%% ---------------------------------------------------------------

\newif\if@inexe@firstline

\newenvironment{inexe}{%
  \inexe@savedleft=\@totalleftmargin
  \advance\inexe@savedleft\leftskip
  \ifnum\value{inexecnt}<9
    \settowidth{\inexe@labelwidth}{\inexe@fmt{9}}%
  \else\ifnum\value{inexecnt}<99
    \settowidth{\inexe@labelwidth}{\inexe@fmt{99}}%
  \else
    \settowidth{\inexe@labelwidth}{\inexe@fmt{999}}%
  \fi\fi
  \@inexe@firstlinetrue
  \leavevmode\newline\noindent
  \let\ex\inexe@ex
  \let\exj\inexe@exj
  \let\sn\inexe@sn
  \let\snj\inexe@snj
  \let\exi\inexe@exi
}{%
  \vspace{-\inexe@itemsep}%
  \vspace{\belowdisplayskip}%
  \ignorespacesafterend
}

\makeatother

\endinput
 in your preamble.
%%
%% Configuration (set these BEFORE %% inexe.tex
%% A display-math-style linguistic example environment.
%% Load with \input{inexe} in your preamble.
%%
%% Configuration (set these BEFORE \input{inexe}):
%%
%%   \def\inexedelimiter{parens}   % or: brackets, period
%%   \newlength{\inexelabelsep}
%%   \setlength{\inexelabelsep}{1em}
%%
%% Counter sharing with gb4e:
%%   If gb4e is loaded, inexe automatically shares gb4e's \exx counter.
%%
%% Usage:
%%   \begin{inexe}
%%     \ex{Consider sub-examples:
%%       \begin{inxlist}     % a. b. c.  (default, like xlist)
%%         \xi First sub-example.
%%         \xij[*] Bad sub-example.
%%         \xisn Unnumbered sub-example.
%%       \end{inxlist}}
%%     \ex{Roman sub-examples:
%%       \begin{inxlisti}    % i. ii. iii.  (like xlisti)
%%         \xi First.
%%         \xi Second.
%%       \end{inxlisti}}
%%   \end{inexe}
%%
%% Inxlist variants (all support \xi, \xij, \xisn, \xisnj):
%%   inxlist   a. b. c.      (mirrors xlist/xlista)
%%   inxlisti  i. ii. iii.   (mirrors xlisti)
%%   inxlistn  1. 2. 3.      (mirrors xlistn)
%%   inxlistA  A. B. C.      (mirrors xlistA)
%%   inxlistI  I. II. III.   (mirrors xlistI)

%% ---------------------------------------------------------------
%% Guard against double-loading
%% ---------------------------------------------------------------
\expandafter\ifx\csname inexe@loaded\endcsname\relax
\expandafter\def\csname inexe@loaded\endcsname{}
\else
  \expandafter\endinput
\fi

\makeatletter

%% ---------------------------------------------------------------
%% User-configurable defaults
%% (Only set if not already defined by the user before \input.)
%% ---------------------------------------------------------------

\@ifundefined{inexedelimiter}{\def\inexedelimiter{parens}}{}

\@ifundefined{inexelabelsep}{%
  \newlength{\inexelabelsep}%
  \setlength{\inexelabelsep}{1em}%
}{}

%% ---------------------------------------------------------------
%% Label formatting
%% ---------------------------------------------------------------

\def\inexe@fmt#1{%
  \ifx\inexedelimiter\@inexe@parens (#1)\else
  \ifx\inexedelimiter\@inexe@brackets [#1]\else
  \ifx\inexedelimiter\@inexe@period #1.\else
  (#1)%
  \fi\fi\fi
}
\def\@inexe@parens{parens}
\def\@inexe@brackets{brackets}
\def\@inexe@period{period}

%% ---------------------------------------------------------------
%% Counter
%% ---------------------------------------------------------------

\@definecounter{inexecnt}

\def\inexe@stepcounter{%
  \refstepcounter{inexecnt}%
  \protected@edef\@currentlabel{\arabic{inexecnt}}%
}
\def\inexe@thecounter{\arabic{inexecnt}}

\AtBeginDocument{%
  \@ifundefined{c@exx}{}{%
    \def\inexe@stepcounter{%
      \refstepcounter{exx}%
      \setcounter{inexecnt}{\value{exx}}%
      \protected@edef\@currentlabel{\arabic{exx}}%
    }%
    \def\inexe@thecounter{\arabic{exx}}%
    \def\theinexecnt{\arabic{exx}}%
  }%
}

%% ---------------------------------------------------------------
%% Core display line
%% ---------------------------------------------------------------

\newlength{\inexe@savedleft}
\newlength{\inexe@labelwidth}
\newlength{\inexe@itemsep}
\setlength{\inexe@itemsep}{3\p@ plus2\p@}

\long\def\inexe@line#1#2{%
  \if@inexe@firstline
    \@inexe@firstlinefalse
  \else
    \par
  \fi
  \noindent
  \hspace*{-\inexe@savedleft}%
  \noindent
  \hbox to \dimexpr\linewidth+\inexe@savedleft\relax{%
    \hbox to \inexe@labelwidth{%
      \hfil
      \ifx\relax#1\relax
      \else
        \inexe@fmt{#1}%
      \fi
    }%
    \hspace{\inexelabelsep}%
    \parbox[t]{\dimexpr
        \linewidth + \inexe@savedleft
        - \inexe@labelwidth
        - \inexelabelsep
      \relax}{#2}%
    \hss
  }%
  \par
  \vspace{\inexe@itemsep}%
}

%% ---------------------------------------------------------------
%% Item commands
%% ---------------------------------------------------------------

\long\def\inexe@ex#1{%
  \inexe@stepcounter
  \inexe@line{\inexe@thecounter}{#1}%
}

\long\def\inexe@exj[#1]#2{%
  \inexe@stepcounter
  \inexe@line{\inexe@thecounter}{#1\hspace{.5ex}#2}%
}

\long\def\inexe@sn#1{%
  \inexe@line{\relax}{#1}%
}

\long\def\inexe@snj[#1]#2{%
  \inexe@line{\relax}{#1\hspace{.5ex}#2}%
}

\long\def\inexe@exi#1#2{%
  \inexe@line{#1}{#2}%
}

%% ---------------------------------------------------------------
%% inxlist — inline sub-example list for use inside \ex{...}
%%
%% Lives inside the \parbox of \ex{...}. Uses \halign so \xi takes
%% unbraced content (everything to the next \xi/\xij/\xisn/\xisnj
%% or \end{inxlist}), matching gb4e's xlist feel.
%% Labels are always a. b. c. style (gb4e xlist default).
%%
%% Usage:
%%   \ex{Consider the following:
%%     \begin{inxlist}
%%       \xi First sub-example.
%%       \xi Second sub-example.
%%       \xij[*] Bad sub-example.
%%       \xisn Unnumbered sub-example.
%%     \end{inxlist}}
%%
%% Commands inside inxlist:
%%   \xi          lettered sub-example (unbraced)
%%   \xij[jdg]    with judgement mark
%%   \xisn        unnumbered sub-example
%%   \xisnj[jdg]  unnumbered with judgement
%% ---------------------------------------------------------------

\@definecounter{inxcnt}
% \p@inxcnt: prefix for \ref — produces "3b" style cross-references.
% Defined as the current inexe label (arabic) so sub-example labels
% inherit the parent example number.
\def\p@inxcnt{\theinexecnt}
% \theinexecnt: the current top-level example number, used as prefix.
% Points to exx if gb4e loaded, else inexecnt. Defined at \begin{document}
% (after gb4e has loaded) via \AtBeginDocument below — placeholder for now.
\def\theinexecnt{\arabic{inexecnt}}
% Save the standard \@item before any package (e.g. paralist) redefines it.
\let\@orignalitem\@item

\newlength{\inxlist@labelwidth}
\newlength{\inxlist@labelsep}
\setlength{\inxlist@labelsep}{.75em}
\newlength{\inxlist@itemsep}
\setlength{\inxlist@itemsep}{1.5\p@ plus\p@}

% \theinxcnt: formats the sub-example counter value (letter/roman/etc.)
% Set locally by \inxlist@open to match the current list style.
\def\theinxcnt{\alph{inxcnt}}
%% Shared setup for all inxlist variants.
%% #1 = counter format command, e.g. \alph, \roman, \arabic, \Alph, \Roman
%% #2 = sample string for label width, e.g. "m." or "iii." or "99."
\def\inxlist@open#1#2{%
  \setcounter{inxcnt}{0}%
  \def\inxlist@counterfmt{#1}%
  % Set \theinxcnt so \refstepcounter{inxcnt} records the right label.
  \def\theinxcnt{#1{inxcnt}}%
  \settowidth{\inxlist@labelwidth}{#2}%
  \edef\inxlist@savedlistdepth{\the\@listdepth}%
  \edef\inxlist@saveditemdepth{\the\@itemdepth}%
  \let\@item\@orignalitem
  \begin{list}{}{%
    \labelwidth=\inxlist@labelwidth
    \labelsep=\inxlist@labelsep
    \leftmargin=\inxlist@labelwidth
    \advance\leftmargin\inxlist@labelsep
    \itemindent\z@
    \listparindent\z@
    \topsep=\inxlist@itemsep
    \itemsep=\inxlist@itemsep
    \parsep\z@
  }%
  \def\xi{%
    \refstepcounter{inxcnt}%
    \item[\inxlist@counterfmt{inxcnt}.]%
  }%
  \def\xij[##1]{%
    \refstepcounter{inxcnt}%
    \item[\inxlist@counterfmt{inxcnt}.]##1\hspace{.5ex}%
  }%
  \def\xisn{\item[]}%
  \def\xisnj[##1]{\item[]##1\hspace{.5ex}}%
}

\def\inxlist@close{%
  \end{list}%
  \@listdepth=\inxlist@savedlistdepth\relax
  \@itemdepth=\inxlist@saveditemdepth\relax
  \ignorespacesafterend
}

% inxlist  — alphabetic:  a. b. c.   (mirrors gb4e's xlist / xlista)
\newenvironment{inxlist}
  {\inxlist@open{\alph}{m.}}
  {\inxlist@close}

% inxlisti — roman:       i. ii. iii.  (mirrors gb4e's xlisti)
\newenvironment{inxlisti}
  {\inxlist@open{\roman}{iii.}}
  {\inxlist@close}

% inxlistn — arabic:      1. 2. 3.   (mirrors gb4e's xlistn)
\newenvironment{inxlistn}
  {\inxlist@open{\arabic}{9.}}
  {\inxlist@close}

% inxlistA — Alphabetic:  A. B. C.   (mirrors gb4e's xlistA)
\newenvironment{inxlistA}
  {\inxlist@open{\Alph}{M.}}
  {\inxlist@close}

% inxlistI — Roman:       I. II. III.  (mirrors gb4e's xlistI)
\newenvironment{inxlistI}
  {\inxlist@open{\Roman}{III.}}
  {\inxlist@close}

%% ---------------------------------------------------------------
%% inexe environment
%% ---------------------------------------------------------------

\newif\if@inexe@firstline

\newenvironment{inexe}{%
  \inexe@savedleft=\@totalleftmargin
  \advance\inexe@savedleft\leftskip
  \ifnum\value{inexecnt}<9
    \settowidth{\inexe@labelwidth}{\inexe@fmt{9}}%
  \else\ifnum\value{inexecnt}<99
    \settowidth{\inexe@labelwidth}{\inexe@fmt{99}}%
  \else
    \settowidth{\inexe@labelwidth}{\inexe@fmt{999}}%
  \fi\fi
  \@inexe@firstlinetrue
  \leavevmode\newline\noindent
  \let\ex\inexe@ex
  \let\exj\inexe@exj
  \let\sn\inexe@sn
  \let\snj\inexe@snj
  \let\exi\inexe@exi
}{%
  \vspace{-\inexe@itemsep}%
  \vspace{\belowdisplayskip}%
  \ignorespacesafterend
}

\makeatother

\endinput
):
%%
%%   \def\inexedelimiter{parens}   % or: brackets, period
%%   \newlength{\inexelabelsep}
%%   \setlength{\inexelabelsep}{1em}
%%
%% Counter sharing with gb4e:
%%   If gb4e is loaded, inexe automatically shares gb4e's \exx counter.
%%
%% Usage:
%%   \begin{inexe}
%%     \ex{Consider sub-examples:
%%       \begin{inxlist}     % a. b. c.  (default, like xlist)
%%         \xi First sub-example.
%%         \xij[*] Bad sub-example.
%%         \xisn Unnumbered sub-example.
%%       \end{inxlist}}
%%     \ex{Roman sub-examples:
%%       \begin{inxlisti}    % i. ii. iii.  (like xlisti)
%%         \xi First.
%%         \xi Second.
%%       \end{inxlisti}}
%%   \end{inexe}
%%
%% Inxlist variants (all support \xi, \xij, \xisn, \xisnj):
%%   inxlist   a. b. c.      (mirrors xlist/xlista)
%%   inxlisti  i. ii. iii.   (mirrors xlisti)
%%   inxlistn  1. 2. 3.      (mirrors xlistn)
%%   inxlistA  A. B. C.      (mirrors xlistA)
%%   inxlistI  I. II. III.   (mirrors xlistI)

%% ---------------------------------------------------------------
%% Guard against double-loading
%% ---------------------------------------------------------------
\expandafter\ifx\csname inexe@loaded\endcsname\relax
\expandafter\def\csname inexe@loaded\endcsname{}
\else
  \expandafter\endinput
\fi

\makeatletter

%% ---------------------------------------------------------------
%% User-configurable defaults
%% (Only set if not already defined by the user before \input.)
%% ---------------------------------------------------------------

\@ifundefined{inexedelimiter}{\def\inexedelimiter{parens}}{}

\@ifundefined{inexelabelsep}{%
  \newlength{\inexelabelsep}%
  \setlength{\inexelabelsep}{1em}%
}{}

%% ---------------------------------------------------------------
%% Label formatting
%% ---------------------------------------------------------------

\def\inexe@fmt#1{%
  \ifx\inexedelimiter\@inexe@parens (#1)\else
  \ifx\inexedelimiter\@inexe@brackets [#1]\else
  \ifx\inexedelimiter\@inexe@period #1.\else
  (#1)%
  \fi\fi\fi
}
\def\@inexe@parens{parens}
\def\@inexe@brackets{brackets}
\def\@inexe@period{period}

%% ---------------------------------------------------------------
%% Counter
%% ---------------------------------------------------------------

\@definecounter{inexecnt}

\def\inexe@stepcounter{%
  \refstepcounter{inexecnt}%
  \protected@edef\@currentlabel{\arabic{inexecnt}}%
}
\def\inexe@thecounter{\arabic{inexecnt}}

\AtBeginDocument{%
  \@ifundefined{c@exx}{}{%
    \def\inexe@stepcounter{%
      \refstepcounter{exx}%
      \setcounter{inexecnt}{\value{exx}}%
      \protected@edef\@currentlabel{\arabic{exx}}%
    }%
    \def\inexe@thecounter{\arabic{exx}}%
    \def\theinexecnt{\arabic{exx}}%
  }%
}

%% ---------------------------------------------------------------
%% Core display line
%% ---------------------------------------------------------------

\newlength{\inexe@savedleft}
\newlength{\inexe@labelwidth}
\newlength{\inexe@itemsep}
\setlength{\inexe@itemsep}{3\p@ plus2\p@}

\long\def\inexe@line#1#2{%
  \if@inexe@firstline
    \@inexe@firstlinefalse
  \else
    \par
  \fi
  \noindent
  \hspace*{-\inexe@savedleft}%
  \noindent
  \hbox to \dimexpr\linewidth+\inexe@savedleft\relax{%
    \hbox to \inexe@labelwidth{%
      \hfil
      \ifx\relax#1\relax
      \else
        \inexe@fmt{#1}%
      \fi
    }%
    \hspace{\inexelabelsep}%
    \parbox[t]{\dimexpr
        \linewidth + \inexe@savedleft
        - \inexe@labelwidth
        - \inexelabelsep
      \relax}{#2}%
    \hss
  }%
  \par
  \vspace{\inexe@itemsep}%
}

%% ---------------------------------------------------------------
%% Item commands
%% ---------------------------------------------------------------

\long\def\inexe@ex#1{%
  \inexe@stepcounter
  \inexe@line{\inexe@thecounter}{#1}%
}

\long\def\inexe@exj[#1]#2{%
  \inexe@stepcounter
  \inexe@line{\inexe@thecounter}{#1\hspace{.5ex}#2}%
}

\long\def\inexe@sn#1{%
  \inexe@line{\relax}{#1}%
}

\long\def\inexe@snj[#1]#2{%
  \inexe@line{\relax}{#1\hspace{.5ex}#2}%
}

\long\def\inexe@exi#1#2{%
  \inexe@line{#1}{#2}%
}

%% ---------------------------------------------------------------
%% inxlist — inline sub-example list for use inside \ex{...}
%%
%% Lives inside the \parbox of \ex{...}. Uses \halign so \xi takes
%% unbraced content (everything to the next \xi/\xij/\xisn/\xisnj
%% or \end{inxlist}), matching gb4e's xlist feel.
%% Labels are always a. b. c. style (gb4e xlist default).
%%
%% Usage:
%%   \ex{Consider the following:
%%     \begin{inxlist}
%%       \xi First sub-example.
%%       \xi Second sub-example.
%%       \xij[*] Bad sub-example.
%%       \xisn Unnumbered sub-example.
%%     \end{inxlist}}
%%
%% Commands inside inxlist:
%%   \xi          lettered sub-example (unbraced)
%%   \xij[jdg]    with judgement mark
%%   \xisn        unnumbered sub-example
%%   \xisnj[jdg]  unnumbered with judgement
%% ---------------------------------------------------------------

\@definecounter{inxcnt}
% \p@inxcnt: prefix for \ref — produces "3b" style cross-references.
% Defined as the current inexe label (arabic) so sub-example labels
% inherit the parent example number.
\def\p@inxcnt{\theinexecnt}
% \theinexecnt: the current top-level example number, used as prefix.
% Points to exx if gb4e loaded, else inexecnt. Defined at \begin{document}
% (after gb4e has loaded) via \AtBeginDocument below — placeholder for now.
\def\theinexecnt{\arabic{inexecnt}}
% Save the standard \@item before any package (e.g. paralist) redefines it.
\let\@orignalitem\@item

\newlength{\inxlist@labelwidth}
\newlength{\inxlist@labelsep}
\setlength{\inxlist@labelsep}{.75em}
\newlength{\inxlist@itemsep}
\setlength{\inxlist@itemsep}{1.5\p@ plus\p@}

% \theinxcnt: formats the sub-example counter value (letter/roman/etc.)
% Set locally by \inxlist@open to match the current list style.
\def\theinxcnt{\alph{inxcnt}}
%% Shared setup for all inxlist variants.
%% #1 = counter format command, e.g. \alph, \roman, \arabic, \Alph, \Roman
%% #2 = sample string for label width, e.g. "m." or "iii." or "99."
\def\inxlist@open#1#2{%
  \setcounter{inxcnt}{0}%
  \def\inxlist@counterfmt{#1}%
  % Set \theinxcnt so \refstepcounter{inxcnt} records the right label.
  \def\theinxcnt{#1{inxcnt}}%
  \settowidth{\inxlist@labelwidth}{#2}%
  \edef\inxlist@savedlistdepth{\the\@listdepth}%
  \edef\inxlist@saveditemdepth{\the\@itemdepth}%
  \let\@item\@orignalitem
  \begin{list}{}{%
    \labelwidth=\inxlist@labelwidth
    \labelsep=\inxlist@labelsep
    \leftmargin=\inxlist@labelwidth
    \advance\leftmargin\inxlist@labelsep
    \itemindent\z@
    \listparindent\z@
    \topsep=\inxlist@itemsep
    \itemsep=\inxlist@itemsep
    \parsep\z@
  }%
  \def\xi{%
    \refstepcounter{inxcnt}%
    \item[\inxlist@counterfmt{inxcnt}.]%
  }%
  \def\xij[##1]{%
    \refstepcounter{inxcnt}%
    \item[\inxlist@counterfmt{inxcnt}.]##1\hspace{.5ex}%
  }%
  \def\xisn{\item[]}%
  \def\xisnj[##1]{\item[]##1\hspace{.5ex}}%
}

\def\inxlist@close{%
  \end{list}%
  \@listdepth=\inxlist@savedlistdepth\relax
  \@itemdepth=\inxlist@saveditemdepth\relax
  \ignorespacesafterend
}

% inxlist  — alphabetic:  a. b. c.   (mirrors gb4e's xlist / xlista)
\newenvironment{inxlist}
  {\inxlist@open{\alph}{m.}}
  {\inxlist@close}

% inxlisti — roman:       i. ii. iii.  (mirrors gb4e's xlisti)
\newenvironment{inxlisti}
  {\inxlist@open{\roman}{iii.}}
  {\inxlist@close}

% inxlistn — arabic:      1. 2. 3.   (mirrors gb4e's xlistn)
\newenvironment{inxlistn}
  {\inxlist@open{\arabic}{9.}}
  {\inxlist@close}

% inxlistA — Alphabetic:  A. B. C.   (mirrors gb4e's xlistA)
\newenvironment{inxlistA}
  {\inxlist@open{\Alph}{M.}}
  {\inxlist@close}

% inxlistI — Roman:       I. II. III.  (mirrors gb4e's xlistI)
\newenvironment{inxlistI}
  {\inxlist@open{\Roman}{III.}}
  {\inxlist@close}

%% ---------------------------------------------------------------
%% inexe environment
%% ---------------------------------------------------------------

\newif\if@inexe@firstline

\newenvironment{inexe}{%
  \inexe@savedleft=\@totalleftmargin
  \advance\inexe@savedleft\leftskip
  \ifnum\value{inexecnt}<9
    \settowidth{\inexe@labelwidth}{\inexe@fmt{9}}%
  \else\ifnum\value{inexecnt}<99
    \settowidth{\inexe@labelwidth}{\inexe@fmt{99}}%
  \else
    \settowidth{\inexe@labelwidth}{\inexe@fmt{999}}%
  \fi\fi
  \@inexe@firstlinetrue
  \leavevmode\newline\noindent
  \let\ex\inexe@ex
  \let\exj\inexe@exj
  \let\sn\inexe@sn
  \let\snj\inexe@snj
  \let\exi\inexe@exi
}{%
  \vspace{-\inexe@itemsep}%
  \vspace{\belowdisplayskip}%
  \ignorespacesafterend
}

\makeatother

\endinput
):
%%
%%   \def\inexedelimiter{parens}   % or: brackets, period
%%   \newlength{\inexelabelsep}
%%   \setlength{\inexelabelsep}{1em}
%%
%% Counter sharing with gb4e:
%%   If gb4e is loaded, inexe automatically shares gb4e's \exx counter.
%%
%% Usage:
%%   \begin{inexe}
%%     \ex{Consider sub-examples:
%%       \begin{inxlist}     % a. b. c.  (default, like xlist)
%%         \xi First sub-example.
%%         \xij[*] Bad sub-example.
%%         \xisn Unnumbered sub-example.
%%       \end{inxlist}}
%%     \ex{Roman sub-examples:
%%       \begin{inxlisti}    % i. ii. iii.  (like xlisti)
%%         \xi First.
%%         \xi Second.
%%       \end{inxlisti}}
%%   \end{inexe}
%%
%% Inxlist variants (all support \xi, \xij, \xisn, \xisnj):
%%   inxlist   a. b. c.      (mirrors xlist/xlista)
%%   inxlisti  i. ii. iii.   (mirrors xlisti)
%%   inxlistn  1. 2. 3.      (mirrors xlistn)
%%   inxlistA  A. B. C.      (mirrors xlistA)
%%   inxlistI  I. II. III.   (mirrors xlistI)

%% ---------------------------------------------------------------
%% Guard against double-loading
%% ---------------------------------------------------------------
\expandafter\ifx\csname inexe@loaded\endcsname\relax
\expandafter\def\csname inexe@loaded\endcsname{}
\else
  \expandafter\endinput
\fi

\makeatletter

%% ---------------------------------------------------------------
%% User-configurable defaults
%% (Only set if not already defined by the user before \input.)
%% ---------------------------------------------------------------

\@ifundefined{inexedelimiter}{\def\inexedelimiter{parens}}{}

\@ifundefined{inexelabelsep}{%
  \newlength{\inexelabelsep}%
  \setlength{\inexelabelsep}{1em}%
}{}

%% ---------------------------------------------------------------
%% Label formatting
%% ---------------------------------------------------------------

\def\inexe@fmt#1{%
  \ifx\inexedelimiter\@inexe@parens (#1)\else
  \ifx\inexedelimiter\@inexe@brackets [#1]\else
  \ifx\inexedelimiter\@inexe@period #1.\else
  (#1)%
  \fi\fi\fi
}
\def\@inexe@parens{parens}
\def\@inexe@brackets{brackets}
\def\@inexe@period{period}

%% ---------------------------------------------------------------
%% Counter
%% ---------------------------------------------------------------

\@definecounter{inexecnt}

\def\inexe@stepcounter{%
  \refstepcounter{inexecnt}%
  \protected@edef\@currentlabel{\arabic{inexecnt}}%
}
\def\inexe@thecounter{\arabic{inexecnt}}

\AtBeginDocument{%
  \@ifundefined{c@exx}{}{%
    \def\inexe@stepcounter{%
      \refstepcounter{exx}%
      \setcounter{inexecnt}{\value{exx}}%
      \protected@edef\@currentlabel{\arabic{exx}}%
    }%
    \def\inexe@thecounter{\arabic{exx}}%
    \def\theinexecnt{\arabic{exx}}%
  }%
}

%% ---------------------------------------------------------------
%% Core display line
%% ---------------------------------------------------------------

\newlength{\inexe@savedleft}
\newlength{\inexe@labelwidth}
\newlength{\inexe@itemsep}
\setlength{\inexe@itemsep}{3\p@ plus2\p@}

\long\def\inexe@line#1#2{%
  \if@inexe@firstline
    \@inexe@firstlinefalse
  \else
    \par
  \fi
  \noindent
  \hspace*{-\inexe@savedleft}%
  \noindent
  \hbox to \dimexpr\linewidth+\inexe@savedleft\relax{%
    \hbox to \inexe@labelwidth{%
      \hfil
      \ifx\relax#1\relax
      \else
        \inexe@fmt{#1}%
      \fi
    }%
    \hspace{\inexelabelsep}%
    \parbox[t]{\dimexpr
        \linewidth + \inexe@savedleft
        - \inexe@labelwidth
        - \inexelabelsep
      \relax}{#2}%
    \hss
  }%
  \par
  \vspace{\inexe@itemsep}%
}

%% ---------------------------------------------------------------
%% Item commands
%% ---------------------------------------------------------------

\long\def\inexe@ex#1{%
  \inexe@stepcounter
  \inexe@line{\inexe@thecounter}{#1}%
}

\long\def\inexe@exj[#1]#2{%
  \inexe@stepcounter
  \inexe@line{\inexe@thecounter}{#1\hspace{.5ex}#2}%
}

\long\def\inexe@sn#1{%
  \inexe@line{\relax}{#1}%
}

\long\def\inexe@snj[#1]#2{%
  \inexe@line{\relax}{#1\hspace{.5ex}#2}%
}

\long\def\inexe@exi#1#2{%
  \inexe@line{#1}{#2}%
}

%% ---------------------------------------------------------------
%% inxlist — inline sub-example list for use inside \ex{...}
%%
%% Lives inside the \parbox of \ex{...}. Uses \halign so \xi takes
%% unbraced content (everything to the next \xi/\xij/\xisn/\xisnj
%% or \end{inxlist}), matching gb4e's xlist feel.
%% Labels are always a. b. c. style (gb4e xlist default).
%%
%% Usage:
%%   \ex{Consider the following:
%%     \begin{inxlist}
%%       \xi First sub-example.
%%       \xi Second sub-example.
%%       \xij[*] Bad sub-example.
%%       \xisn Unnumbered sub-example.
%%     \end{inxlist}}
%%
%% Commands inside inxlist:
%%   \xi          lettered sub-example (unbraced)
%%   \xij[jdg]    with judgement mark
%%   \xisn        unnumbered sub-example
%%   \xisnj[jdg]  unnumbered with judgement
%% ---------------------------------------------------------------

\@definecounter{inxcnt}
% \p@inxcnt: prefix for \ref — produces "3b" style cross-references.
% Defined as the current inexe label (arabic) so sub-example labels
% inherit the parent example number.
\def\p@inxcnt{\theinexecnt}
% \theinexecnt: the current top-level example number, used as prefix.
% Points to exx if gb4e loaded, else inexecnt. Defined at \begin{document}
% (after gb4e has loaded) via \AtBeginDocument below — placeholder for now.
\def\theinexecnt{\arabic{inexecnt}}
% Save the standard \@item before any package (e.g. paralist) redefines it.
\let\@orignalitem\@item

\newlength{\inxlist@labelwidth}
\newlength{\inxlist@labelsep}
\setlength{\inxlist@labelsep}{.75em}
\newlength{\inxlist@itemsep}
\setlength{\inxlist@itemsep}{1.5\p@ plus\p@}

% \theinxcnt: formats the sub-example counter value (letter/roman/etc.)
% Set locally by \inxlist@open to match the current list style.
\def\theinxcnt{\alph{inxcnt}}
%% Shared setup for all inxlist variants.
%% #1 = counter format command, e.g. \alph, \roman, \arabic, \Alph, \Roman
%% #2 = sample string for label width, e.g. "m." or "iii." or "99."
\def\inxlist@open#1#2{%
  \setcounter{inxcnt}{0}%
  \def\inxlist@counterfmt{#1}%
  % Set \theinxcnt so \refstepcounter{inxcnt} records the right label.
  \def\theinxcnt{#1{inxcnt}}%
  \settowidth{\inxlist@labelwidth}{#2}%
  \edef\inxlist@savedlistdepth{\the\@listdepth}%
  \edef\inxlist@saveditemdepth{\the\@itemdepth}%
  \let\@item\@orignalitem
  \begin{list}{}{%
    \labelwidth=\inxlist@labelwidth
    \labelsep=\inxlist@labelsep
    \leftmargin=\inxlist@labelwidth
    \advance\leftmargin\inxlist@labelsep
    \itemindent\z@
    \listparindent\z@
    \topsep=\inxlist@itemsep
    \itemsep=\inxlist@itemsep
    \parsep\z@
  }%
  \def\xi{%
    \refstepcounter{inxcnt}%
    \item[\inxlist@counterfmt{inxcnt}.]%
  }%
  \def\xij[##1]{%
    \refstepcounter{inxcnt}%
    \item[\inxlist@counterfmt{inxcnt}.]##1\hspace{.5ex}%
  }%
  \def\xisn{\item[]}%
  \def\xisnj[##1]{\item[]##1\hspace{.5ex}}%
}

\def\inxlist@close{%
  \end{list}%
  \@listdepth=\inxlist@savedlistdepth\relax
  \@itemdepth=\inxlist@saveditemdepth\relax
  \ignorespacesafterend
}

% inxlist  — alphabetic:  a. b. c.   (mirrors gb4e's xlist / xlista)
\newenvironment{inxlist}
  {\inxlist@open{\alph}{m.}}
  {\inxlist@close}

% inxlisti — roman:       i. ii. iii.  (mirrors gb4e's xlisti)
\newenvironment{inxlisti}
  {\inxlist@open{\roman}{iii.}}
  {\inxlist@close}

% inxlistn — arabic:      1. 2. 3.   (mirrors gb4e's xlistn)
\newenvironment{inxlistn}
  {\inxlist@open{\arabic}{9.}}
  {\inxlist@close}

% inxlistA — Alphabetic:  A. B. C.   (mirrors gb4e's xlistA)
\newenvironment{inxlistA}
  {\inxlist@open{\Alph}{M.}}
  {\inxlist@close}

% inxlistI — Roman:       I. II. III.  (mirrors gb4e's xlistI)
\newenvironment{inxlistI}
  {\inxlist@open{\Roman}{III.}}
  {\inxlist@close}

%% ---------------------------------------------------------------
%% inexe environment
%% ---------------------------------------------------------------

\newif\if@inexe@firstline

\newenvironment{inexe}{%
  \inexe@savedleft=\@totalleftmargin
  \advance\inexe@savedleft\leftskip
  \ifnum\value{inexecnt}<9
    \settowidth{\inexe@labelwidth}{\inexe@fmt{9}}%
  \else\ifnum\value{inexecnt}<99
    \settowidth{\inexe@labelwidth}{\inexe@fmt{99}}%
  \else
    \settowidth{\inexe@labelwidth}{\inexe@fmt{999}}%
  \fi\fi
  \@inexe@firstlinetrue
  \leavevmode\newline\noindent
  \let\ex\inexe@ex
  \let\exj\inexe@exj
  \let\sn\inexe@sn
  \let\snj\inexe@snj
  \let\exi\inexe@exi
}{%
  \vspace{-\inexe@itemsep}%
  \vspace{\belowdisplayskip}%
  \ignorespacesafterend
}

\makeatother

\endinput
 in your preamble.
%%
%% Configuration (set these BEFORE %% inexe.tex
%% A display-math-style linguistic example environment.
%% Load with %% inexe.tex
%% A display-math-style linguistic example environment.
%% Load with %% inexe.tex
%% A display-math-style linguistic example environment.
%% Load with \input{inexe} in your preamble.
%%
%% Configuration (set these BEFORE \input{inexe}):
%%
%%   \def\inexedelimiter{parens}   % or: brackets, period
%%   \newlength{\inexelabelsep}
%%   \setlength{\inexelabelsep}{1em}
%%
%% Counter sharing with gb4e:
%%   If gb4e is loaded, inexe automatically shares gb4e's \exx counter.
%%
%% Usage:
%%   \begin{inexe}
%%     \ex{Consider sub-examples:
%%       \begin{inxlist}     % a. b. c.  (default, like xlist)
%%         \xi First sub-example.
%%         \xij[*] Bad sub-example.
%%         \xisn Unnumbered sub-example.
%%       \end{inxlist}}
%%     \ex{Roman sub-examples:
%%       \begin{inxlisti}    % i. ii. iii.  (like xlisti)
%%         \xi First.
%%         \xi Second.
%%       \end{inxlisti}}
%%   \end{inexe}
%%
%% Inxlist variants (all support \xi, \xij, \xisn, \xisnj):
%%   inxlist   a. b. c.      (mirrors xlist/xlista)
%%   inxlisti  i. ii. iii.   (mirrors xlisti)
%%   inxlistn  1. 2. 3.      (mirrors xlistn)
%%   inxlistA  A. B. C.      (mirrors xlistA)
%%   inxlistI  I. II. III.   (mirrors xlistI)

%% ---------------------------------------------------------------
%% Guard against double-loading
%% ---------------------------------------------------------------
\expandafter\ifx\csname inexe@loaded\endcsname\relax
\expandafter\def\csname inexe@loaded\endcsname{}
\else
  \expandafter\endinput
\fi

\makeatletter

%% ---------------------------------------------------------------
%% User-configurable defaults
%% (Only set if not already defined by the user before \input.)
%% ---------------------------------------------------------------

\@ifundefined{inexedelimiter}{\def\inexedelimiter{parens}}{}

\@ifundefined{inexelabelsep}{%
  \newlength{\inexelabelsep}%
  \setlength{\inexelabelsep}{1em}%
}{}

%% ---------------------------------------------------------------
%% Label formatting
%% ---------------------------------------------------------------

\def\inexe@fmt#1{%
  \ifx\inexedelimiter\@inexe@parens (#1)\else
  \ifx\inexedelimiter\@inexe@brackets [#1]\else
  \ifx\inexedelimiter\@inexe@period #1.\else
  (#1)%
  \fi\fi\fi
}
\def\@inexe@parens{parens}
\def\@inexe@brackets{brackets}
\def\@inexe@period{period}

%% ---------------------------------------------------------------
%% Counter
%% ---------------------------------------------------------------

\@definecounter{inexecnt}

\def\inexe@stepcounter{%
  \refstepcounter{inexecnt}%
  \protected@edef\@currentlabel{\arabic{inexecnt}}%
}
\def\inexe@thecounter{\arabic{inexecnt}}

\AtBeginDocument{%
  \@ifundefined{c@exx}{}{%
    \def\inexe@stepcounter{%
      \refstepcounter{exx}%
      \setcounter{inexecnt}{\value{exx}}%
      \protected@edef\@currentlabel{\arabic{exx}}%
    }%
    \def\inexe@thecounter{\arabic{exx}}%
    \def\theinexecnt{\arabic{exx}}%
  }%
}

%% ---------------------------------------------------------------
%% Core display line
%% ---------------------------------------------------------------

\newlength{\inexe@savedleft}
\newlength{\inexe@labelwidth}
\newlength{\inexe@itemsep}
\setlength{\inexe@itemsep}{3\p@ plus2\p@}

\long\def\inexe@line#1#2{%
  \if@inexe@firstline
    \@inexe@firstlinefalse
  \else
    \par
  \fi
  \noindent
  \hspace*{-\inexe@savedleft}%
  \noindent
  \hbox to \dimexpr\linewidth+\inexe@savedleft\relax{%
    \hbox to \inexe@labelwidth{%
      \hfil
      \ifx\relax#1\relax
      \else
        \inexe@fmt{#1}%
      \fi
    }%
    \hspace{\inexelabelsep}%
    \parbox[t]{\dimexpr
        \linewidth + \inexe@savedleft
        - \inexe@labelwidth
        - \inexelabelsep
      \relax}{#2}%
    \hss
  }%
  \par
  \vspace{\inexe@itemsep}%
}

%% ---------------------------------------------------------------
%% Item commands
%% ---------------------------------------------------------------

\long\def\inexe@ex#1{%
  \inexe@stepcounter
  \inexe@line{\inexe@thecounter}{#1}%
}

\long\def\inexe@exj[#1]#2{%
  \inexe@stepcounter
  \inexe@line{\inexe@thecounter}{#1\hspace{.5ex}#2}%
}

\long\def\inexe@sn#1{%
  \inexe@line{\relax}{#1}%
}

\long\def\inexe@snj[#1]#2{%
  \inexe@line{\relax}{#1\hspace{.5ex}#2}%
}

\long\def\inexe@exi#1#2{%
  \inexe@line{#1}{#2}%
}

%% ---------------------------------------------------------------
%% inxlist — inline sub-example list for use inside \ex{...}
%%
%% Lives inside the \parbox of \ex{...}. Uses \halign so \xi takes
%% unbraced content (everything to the next \xi/\xij/\xisn/\xisnj
%% or \end{inxlist}), matching gb4e's xlist feel.
%% Labels are always a. b. c. style (gb4e xlist default).
%%
%% Usage:
%%   \ex{Consider the following:
%%     \begin{inxlist}
%%       \xi First sub-example.
%%       \xi Second sub-example.
%%       \xij[*] Bad sub-example.
%%       \xisn Unnumbered sub-example.
%%     \end{inxlist}}
%%
%% Commands inside inxlist:
%%   \xi          lettered sub-example (unbraced)
%%   \xij[jdg]    with judgement mark
%%   \xisn        unnumbered sub-example
%%   \xisnj[jdg]  unnumbered with judgement
%% ---------------------------------------------------------------

\@definecounter{inxcnt}
% \p@inxcnt: prefix for \ref — produces "3b" style cross-references.
% Defined as the current inexe label (arabic) so sub-example labels
% inherit the parent example number.
\def\p@inxcnt{\theinexecnt}
% \theinexecnt: the current top-level example number, used as prefix.
% Points to exx if gb4e loaded, else inexecnt. Defined at \begin{document}
% (after gb4e has loaded) via \AtBeginDocument below — placeholder for now.
\def\theinexecnt{\arabic{inexecnt}}
% Save the standard \@item before any package (e.g. paralist) redefines it.
\let\@orignalitem\@item

\newlength{\inxlist@labelwidth}
\newlength{\inxlist@labelsep}
\setlength{\inxlist@labelsep}{.75em}
\newlength{\inxlist@itemsep}
\setlength{\inxlist@itemsep}{1.5\p@ plus\p@}

% \theinxcnt: formats the sub-example counter value (letter/roman/etc.)
% Set locally by \inxlist@open to match the current list style.
\def\theinxcnt{\alph{inxcnt}}
%% Shared setup for all inxlist variants.
%% #1 = counter format command, e.g. \alph, \roman, \arabic, \Alph, \Roman
%% #2 = sample string for label width, e.g. "m." or "iii." or "99."
\def\inxlist@open#1#2{%
  \setcounter{inxcnt}{0}%
  \def\inxlist@counterfmt{#1}%
  % Set \theinxcnt so \refstepcounter{inxcnt} records the right label.
  \def\theinxcnt{#1{inxcnt}}%
  \settowidth{\inxlist@labelwidth}{#2}%
  \edef\inxlist@savedlistdepth{\the\@listdepth}%
  \edef\inxlist@saveditemdepth{\the\@itemdepth}%
  \let\@item\@orignalitem
  \begin{list}{}{%
    \labelwidth=\inxlist@labelwidth
    \labelsep=\inxlist@labelsep
    \leftmargin=\inxlist@labelwidth
    \advance\leftmargin\inxlist@labelsep
    \itemindent\z@
    \listparindent\z@
    \topsep=\inxlist@itemsep
    \itemsep=\inxlist@itemsep
    \parsep\z@
  }%
  \def\xi{%
    \refstepcounter{inxcnt}%
    \item[\inxlist@counterfmt{inxcnt}.]%
  }%
  \def\xij[##1]{%
    \refstepcounter{inxcnt}%
    \item[\inxlist@counterfmt{inxcnt}.]##1\hspace{.5ex}%
  }%
  \def\xisn{\item[]}%
  \def\xisnj[##1]{\item[]##1\hspace{.5ex}}%
}

\def\inxlist@close{%
  \end{list}%
  \@listdepth=\inxlist@savedlistdepth\relax
  \@itemdepth=\inxlist@saveditemdepth\relax
  \ignorespacesafterend
}

% inxlist  — alphabetic:  a. b. c.   (mirrors gb4e's xlist / xlista)
\newenvironment{inxlist}
  {\inxlist@open{\alph}{m.}}
  {\inxlist@close}

% inxlisti — roman:       i. ii. iii.  (mirrors gb4e's xlisti)
\newenvironment{inxlisti}
  {\inxlist@open{\roman}{iii.}}
  {\inxlist@close}

% inxlistn — arabic:      1. 2. 3.   (mirrors gb4e's xlistn)
\newenvironment{inxlistn}
  {\inxlist@open{\arabic}{9.}}
  {\inxlist@close}

% inxlistA — Alphabetic:  A. B. C.   (mirrors gb4e's xlistA)
\newenvironment{inxlistA}
  {\inxlist@open{\Alph}{M.}}
  {\inxlist@close}

% inxlistI — Roman:       I. II. III.  (mirrors gb4e's xlistI)
\newenvironment{inxlistI}
  {\inxlist@open{\Roman}{III.}}
  {\inxlist@close}

%% ---------------------------------------------------------------
%% inexe environment
%% ---------------------------------------------------------------

\newif\if@inexe@firstline

\newenvironment{inexe}{%
  \inexe@savedleft=\@totalleftmargin
  \advance\inexe@savedleft\leftskip
  \ifnum\value{inexecnt}<9
    \settowidth{\inexe@labelwidth}{\inexe@fmt{9}}%
  \else\ifnum\value{inexecnt}<99
    \settowidth{\inexe@labelwidth}{\inexe@fmt{99}}%
  \else
    \settowidth{\inexe@labelwidth}{\inexe@fmt{999}}%
  \fi\fi
  \@inexe@firstlinetrue
  \leavevmode\newline\noindent
  \let\ex\inexe@ex
  \let\exj\inexe@exj
  \let\sn\inexe@sn
  \let\snj\inexe@snj
  \let\exi\inexe@exi
}{%
  \vspace{-\inexe@itemsep}%
  \vspace{\belowdisplayskip}%
  \ignorespacesafterend
}

\makeatother

\endinput
 in your preamble.
%%
%% Configuration (set these BEFORE %% inexe.tex
%% A display-math-style linguistic example environment.
%% Load with \input{inexe} in your preamble.
%%
%% Configuration (set these BEFORE \input{inexe}):
%%
%%   \def\inexedelimiter{parens}   % or: brackets, period
%%   \newlength{\inexelabelsep}
%%   \setlength{\inexelabelsep}{1em}
%%
%% Counter sharing with gb4e:
%%   If gb4e is loaded, inexe automatically shares gb4e's \exx counter.
%%
%% Usage:
%%   \begin{inexe}
%%     \ex{Consider sub-examples:
%%       \begin{inxlist}     % a. b. c.  (default, like xlist)
%%         \xi First sub-example.
%%         \xij[*] Bad sub-example.
%%         \xisn Unnumbered sub-example.
%%       \end{inxlist}}
%%     \ex{Roman sub-examples:
%%       \begin{inxlisti}    % i. ii. iii.  (like xlisti)
%%         \xi First.
%%         \xi Second.
%%       \end{inxlisti}}
%%   \end{inexe}
%%
%% Inxlist variants (all support \xi, \xij, \xisn, \xisnj):
%%   inxlist   a. b. c.      (mirrors xlist/xlista)
%%   inxlisti  i. ii. iii.   (mirrors xlisti)
%%   inxlistn  1. 2. 3.      (mirrors xlistn)
%%   inxlistA  A. B. C.      (mirrors xlistA)
%%   inxlistI  I. II. III.   (mirrors xlistI)

%% ---------------------------------------------------------------
%% Guard against double-loading
%% ---------------------------------------------------------------
\expandafter\ifx\csname inexe@loaded\endcsname\relax
\expandafter\def\csname inexe@loaded\endcsname{}
\else
  \expandafter\endinput
\fi

\makeatletter

%% ---------------------------------------------------------------
%% User-configurable defaults
%% (Only set if not already defined by the user before \input.)
%% ---------------------------------------------------------------

\@ifundefined{inexedelimiter}{\def\inexedelimiter{parens}}{}

\@ifundefined{inexelabelsep}{%
  \newlength{\inexelabelsep}%
  \setlength{\inexelabelsep}{1em}%
}{}

%% ---------------------------------------------------------------
%% Label formatting
%% ---------------------------------------------------------------

\def\inexe@fmt#1{%
  \ifx\inexedelimiter\@inexe@parens (#1)\else
  \ifx\inexedelimiter\@inexe@brackets [#1]\else
  \ifx\inexedelimiter\@inexe@period #1.\else
  (#1)%
  \fi\fi\fi
}
\def\@inexe@parens{parens}
\def\@inexe@brackets{brackets}
\def\@inexe@period{period}

%% ---------------------------------------------------------------
%% Counter
%% ---------------------------------------------------------------

\@definecounter{inexecnt}

\def\inexe@stepcounter{%
  \refstepcounter{inexecnt}%
  \protected@edef\@currentlabel{\arabic{inexecnt}}%
}
\def\inexe@thecounter{\arabic{inexecnt}}

\AtBeginDocument{%
  \@ifundefined{c@exx}{}{%
    \def\inexe@stepcounter{%
      \refstepcounter{exx}%
      \setcounter{inexecnt}{\value{exx}}%
      \protected@edef\@currentlabel{\arabic{exx}}%
    }%
    \def\inexe@thecounter{\arabic{exx}}%
    \def\theinexecnt{\arabic{exx}}%
  }%
}

%% ---------------------------------------------------------------
%% Core display line
%% ---------------------------------------------------------------

\newlength{\inexe@savedleft}
\newlength{\inexe@labelwidth}
\newlength{\inexe@itemsep}
\setlength{\inexe@itemsep}{3\p@ plus2\p@}

\long\def\inexe@line#1#2{%
  \if@inexe@firstline
    \@inexe@firstlinefalse
  \else
    \par
  \fi
  \noindent
  \hspace*{-\inexe@savedleft}%
  \noindent
  \hbox to \dimexpr\linewidth+\inexe@savedleft\relax{%
    \hbox to \inexe@labelwidth{%
      \hfil
      \ifx\relax#1\relax
      \else
        \inexe@fmt{#1}%
      \fi
    }%
    \hspace{\inexelabelsep}%
    \parbox[t]{\dimexpr
        \linewidth + \inexe@savedleft
        - \inexe@labelwidth
        - \inexelabelsep
      \relax}{#2}%
    \hss
  }%
  \par
  \vspace{\inexe@itemsep}%
}

%% ---------------------------------------------------------------
%% Item commands
%% ---------------------------------------------------------------

\long\def\inexe@ex#1{%
  \inexe@stepcounter
  \inexe@line{\inexe@thecounter}{#1}%
}

\long\def\inexe@exj[#1]#2{%
  \inexe@stepcounter
  \inexe@line{\inexe@thecounter}{#1\hspace{.5ex}#2}%
}

\long\def\inexe@sn#1{%
  \inexe@line{\relax}{#1}%
}

\long\def\inexe@snj[#1]#2{%
  \inexe@line{\relax}{#1\hspace{.5ex}#2}%
}

\long\def\inexe@exi#1#2{%
  \inexe@line{#1}{#2}%
}

%% ---------------------------------------------------------------
%% inxlist — inline sub-example list for use inside \ex{...}
%%
%% Lives inside the \parbox of \ex{...}. Uses \halign so \xi takes
%% unbraced content (everything to the next \xi/\xij/\xisn/\xisnj
%% or \end{inxlist}), matching gb4e's xlist feel.
%% Labels are always a. b. c. style (gb4e xlist default).
%%
%% Usage:
%%   \ex{Consider the following:
%%     \begin{inxlist}
%%       \xi First sub-example.
%%       \xi Second sub-example.
%%       \xij[*] Bad sub-example.
%%       \xisn Unnumbered sub-example.
%%     \end{inxlist}}
%%
%% Commands inside inxlist:
%%   \xi          lettered sub-example (unbraced)
%%   \xij[jdg]    with judgement mark
%%   \xisn        unnumbered sub-example
%%   \xisnj[jdg]  unnumbered with judgement
%% ---------------------------------------------------------------

\@definecounter{inxcnt}
% \p@inxcnt: prefix for \ref — produces "3b" style cross-references.
% Defined as the current inexe label (arabic) so sub-example labels
% inherit the parent example number.
\def\p@inxcnt{\theinexecnt}
% \theinexecnt: the current top-level example number, used as prefix.
% Points to exx if gb4e loaded, else inexecnt. Defined at \begin{document}
% (after gb4e has loaded) via \AtBeginDocument below — placeholder for now.
\def\theinexecnt{\arabic{inexecnt}}
% Save the standard \@item before any package (e.g. paralist) redefines it.
\let\@orignalitem\@item

\newlength{\inxlist@labelwidth}
\newlength{\inxlist@labelsep}
\setlength{\inxlist@labelsep}{.75em}
\newlength{\inxlist@itemsep}
\setlength{\inxlist@itemsep}{1.5\p@ plus\p@}

% \theinxcnt: formats the sub-example counter value (letter/roman/etc.)
% Set locally by \inxlist@open to match the current list style.
\def\theinxcnt{\alph{inxcnt}}
%% Shared setup for all inxlist variants.
%% #1 = counter format command, e.g. \alph, \roman, \arabic, \Alph, \Roman
%% #2 = sample string for label width, e.g. "m." or "iii." or "99."
\def\inxlist@open#1#2{%
  \setcounter{inxcnt}{0}%
  \def\inxlist@counterfmt{#1}%
  % Set \theinxcnt so \refstepcounter{inxcnt} records the right label.
  \def\theinxcnt{#1{inxcnt}}%
  \settowidth{\inxlist@labelwidth}{#2}%
  \edef\inxlist@savedlistdepth{\the\@listdepth}%
  \edef\inxlist@saveditemdepth{\the\@itemdepth}%
  \let\@item\@orignalitem
  \begin{list}{}{%
    \labelwidth=\inxlist@labelwidth
    \labelsep=\inxlist@labelsep
    \leftmargin=\inxlist@labelwidth
    \advance\leftmargin\inxlist@labelsep
    \itemindent\z@
    \listparindent\z@
    \topsep=\inxlist@itemsep
    \itemsep=\inxlist@itemsep
    \parsep\z@
  }%
  \def\xi{%
    \refstepcounter{inxcnt}%
    \item[\inxlist@counterfmt{inxcnt}.]%
  }%
  \def\xij[##1]{%
    \refstepcounter{inxcnt}%
    \item[\inxlist@counterfmt{inxcnt}.]##1\hspace{.5ex}%
  }%
  \def\xisn{\item[]}%
  \def\xisnj[##1]{\item[]##1\hspace{.5ex}}%
}

\def\inxlist@close{%
  \end{list}%
  \@listdepth=\inxlist@savedlistdepth\relax
  \@itemdepth=\inxlist@saveditemdepth\relax
  \ignorespacesafterend
}

% inxlist  — alphabetic:  a. b. c.   (mirrors gb4e's xlist / xlista)
\newenvironment{inxlist}
  {\inxlist@open{\alph}{m.}}
  {\inxlist@close}

% inxlisti — roman:       i. ii. iii.  (mirrors gb4e's xlisti)
\newenvironment{inxlisti}
  {\inxlist@open{\roman}{iii.}}
  {\inxlist@close}

% inxlistn — arabic:      1. 2. 3.   (mirrors gb4e's xlistn)
\newenvironment{inxlistn}
  {\inxlist@open{\arabic}{9.}}
  {\inxlist@close}

% inxlistA — Alphabetic:  A. B. C.   (mirrors gb4e's xlistA)
\newenvironment{inxlistA}
  {\inxlist@open{\Alph}{M.}}
  {\inxlist@close}

% inxlistI — Roman:       I. II. III.  (mirrors gb4e's xlistI)
\newenvironment{inxlistI}
  {\inxlist@open{\Roman}{III.}}
  {\inxlist@close}

%% ---------------------------------------------------------------
%% inexe environment
%% ---------------------------------------------------------------

\newif\if@inexe@firstline

\newenvironment{inexe}{%
  \inexe@savedleft=\@totalleftmargin
  \advance\inexe@savedleft\leftskip
  \ifnum\value{inexecnt}<9
    \settowidth{\inexe@labelwidth}{\inexe@fmt{9}}%
  \else\ifnum\value{inexecnt}<99
    \settowidth{\inexe@labelwidth}{\inexe@fmt{99}}%
  \else
    \settowidth{\inexe@labelwidth}{\inexe@fmt{999}}%
  \fi\fi
  \@inexe@firstlinetrue
  \leavevmode\newline\noindent
  \let\ex\inexe@ex
  \let\exj\inexe@exj
  \let\sn\inexe@sn
  \let\snj\inexe@snj
  \let\exi\inexe@exi
}{%
  \vspace{-\inexe@itemsep}%
  \vspace{\belowdisplayskip}%
  \ignorespacesafterend
}

\makeatother

\endinput
):
%%
%%   \def\inexedelimiter{parens}   % or: brackets, period
%%   \newlength{\inexelabelsep}
%%   \setlength{\inexelabelsep}{1em}
%%
%% Counter sharing with gb4e:
%%   If gb4e is loaded, inexe automatically shares gb4e's \exx counter.
%%
%% Usage:
%%   \begin{inexe}
%%     \ex{Consider sub-examples:
%%       \begin{inxlist}     % a. b. c.  (default, like xlist)
%%         \xi First sub-example.
%%         \xij[*] Bad sub-example.
%%         \xisn Unnumbered sub-example.
%%       \end{inxlist}}
%%     \ex{Roman sub-examples:
%%       \begin{inxlisti}    % i. ii. iii.  (like xlisti)
%%         \xi First.
%%         \xi Second.
%%       \end{inxlisti}}
%%   \end{inexe}
%%
%% Inxlist variants (all support \xi, \xij, \xisn, \xisnj):
%%   inxlist   a. b. c.      (mirrors xlist/xlista)
%%   inxlisti  i. ii. iii.   (mirrors xlisti)
%%   inxlistn  1. 2. 3.      (mirrors xlistn)
%%   inxlistA  A. B. C.      (mirrors xlistA)
%%   inxlistI  I. II. III.   (mirrors xlistI)

%% ---------------------------------------------------------------
%% Guard against double-loading
%% ---------------------------------------------------------------
\expandafter\ifx\csname inexe@loaded\endcsname\relax
\expandafter\def\csname inexe@loaded\endcsname{}
\else
  \expandafter\endinput
\fi

\makeatletter

%% ---------------------------------------------------------------
%% User-configurable defaults
%% (Only set if not already defined by the user before \input.)
%% ---------------------------------------------------------------

\@ifundefined{inexedelimiter}{\def\inexedelimiter{parens}}{}

\@ifundefined{inexelabelsep}{%
  \newlength{\inexelabelsep}%
  \setlength{\inexelabelsep}{1em}%
}{}

%% ---------------------------------------------------------------
%% Label formatting
%% ---------------------------------------------------------------

\def\inexe@fmt#1{%
  \ifx\inexedelimiter\@inexe@parens (#1)\else
  \ifx\inexedelimiter\@inexe@brackets [#1]\else
  \ifx\inexedelimiter\@inexe@period #1.\else
  (#1)%
  \fi\fi\fi
}
\def\@inexe@parens{parens}
\def\@inexe@brackets{brackets}
\def\@inexe@period{period}

%% ---------------------------------------------------------------
%% Counter
%% ---------------------------------------------------------------

\@definecounter{inexecnt}

\def\inexe@stepcounter{%
  \refstepcounter{inexecnt}%
  \protected@edef\@currentlabel{\arabic{inexecnt}}%
}
\def\inexe@thecounter{\arabic{inexecnt}}

\AtBeginDocument{%
  \@ifundefined{c@exx}{}{%
    \def\inexe@stepcounter{%
      \refstepcounter{exx}%
      \setcounter{inexecnt}{\value{exx}}%
      \protected@edef\@currentlabel{\arabic{exx}}%
    }%
    \def\inexe@thecounter{\arabic{exx}}%
    \def\theinexecnt{\arabic{exx}}%
  }%
}

%% ---------------------------------------------------------------
%% Core display line
%% ---------------------------------------------------------------

\newlength{\inexe@savedleft}
\newlength{\inexe@labelwidth}
\newlength{\inexe@itemsep}
\setlength{\inexe@itemsep}{3\p@ plus2\p@}

\long\def\inexe@line#1#2{%
  \if@inexe@firstline
    \@inexe@firstlinefalse
  \else
    \par
  \fi
  \noindent
  \hspace*{-\inexe@savedleft}%
  \noindent
  \hbox to \dimexpr\linewidth+\inexe@savedleft\relax{%
    \hbox to \inexe@labelwidth{%
      \hfil
      \ifx\relax#1\relax
      \else
        \inexe@fmt{#1}%
      \fi
    }%
    \hspace{\inexelabelsep}%
    \parbox[t]{\dimexpr
        \linewidth + \inexe@savedleft
        - \inexe@labelwidth
        - \inexelabelsep
      \relax}{#2}%
    \hss
  }%
  \par
  \vspace{\inexe@itemsep}%
}

%% ---------------------------------------------------------------
%% Item commands
%% ---------------------------------------------------------------

\long\def\inexe@ex#1{%
  \inexe@stepcounter
  \inexe@line{\inexe@thecounter}{#1}%
}

\long\def\inexe@exj[#1]#2{%
  \inexe@stepcounter
  \inexe@line{\inexe@thecounter}{#1\hspace{.5ex}#2}%
}

\long\def\inexe@sn#1{%
  \inexe@line{\relax}{#1}%
}

\long\def\inexe@snj[#1]#2{%
  \inexe@line{\relax}{#1\hspace{.5ex}#2}%
}

\long\def\inexe@exi#1#2{%
  \inexe@line{#1}{#2}%
}

%% ---------------------------------------------------------------
%% inxlist — inline sub-example list for use inside \ex{...}
%%
%% Lives inside the \parbox of \ex{...}. Uses \halign so \xi takes
%% unbraced content (everything to the next \xi/\xij/\xisn/\xisnj
%% or \end{inxlist}), matching gb4e's xlist feel.
%% Labels are always a. b. c. style (gb4e xlist default).
%%
%% Usage:
%%   \ex{Consider the following:
%%     \begin{inxlist}
%%       \xi First sub-example.
%%       \xi Second sub-example.
%%       \xij[*] Bad sub-example.
%%       \xisn Unnumbered sub-example.
%%     \end{inxlist}}
%%
%% Commands inside inxlist:
%%   \xi          lettered sub-example (unbraced)
%%   \xij[jdg]    with judgement mark
%%   \xisn        unnumbered sub-example
%%   \xisnj[jdg]  unnumbered with judgement
%% ---------------------------------------------------------------

\@definecounter{inxcnt}
% \p@inxcnt: prefix for \ref — produces "3b" style cross-references.
% Defined as the current inexe label (arabic) so sub-example labels
% inherit the parent example number.
\def\p@inxcnt{\theinexecnt}
% \theinexecnt: the current top-level example number, used as prefix.
% Points to exx if gb4e loaded, else inexecnt. Defined at \begin{document}
% (after gb4e has loaded) via \AtBeginDocument below — placeholder for now.
\def\theinexecnt{\arabic{inexecnt}}
% Save the standard \@item before any package (e.g. paralist) redefines it.
\let\@orignalitem\@item

\newlength{\inxlist@labelwidth}
\newlength{\inxlist@labelsep}
\setlength{\inxlist@labelsep}{.75em}
\newlength{\inxlist@itemsep}
\setlength{\inxlist@itemsep}{1.5\p@ plus\p@}

% \theinxcnt: formats the sub-example counter value (letter/roman/etc.)
% Set locally by \inxlist@open to match the current list style.
\def\theinxcnt{\alph{inxcnt}}
%% Shared setup for all inxlist variants.
%% #1 = counter format command, e.g. \alph, \roman, \arabic, \Alph, \Roman
%% #2 = sample string for label width, e.g. "m." or "iii." or "99."
\def\inxlist@open#1#2{%
  \setcounter{inxcnt}{0}%
  \def\inxlist@counterfmt{#1}%
  % Set \theinxcnt so \refstepcounter{inxcnt} records the right label.
  \def\theinxcnt{#1{inxcnt}}%
  \settowidth{\inxlist@labelwidth}{#2}%
  \edef\inxlist@savedlistdepth{\the\@listdepth}%
  \edef\inxlist@saveditemdepth{\the\@itemdepth}%
  \let\@item\@orignalitem
  \begin{list}{}{%
    \labelwidth=\inxlist@labelwidth
    \labelsep=\inxlist@labelsep
    \leftmargin=\inxlist@labelwidth
    \advance\leftmargin\inxlist@labelsep
    \itemindent\z@
    \listparindent\z@
    \topsep=\inxlist@itemsep
    \itemsep=\inxlist@itemsep
    \parsep\z@
  }%
  \def\xi{%
    \refstepcounter{inxcnt}%
    \item[\inxlist@counterfmt{inxcnt}.]%
  }%
  \def\xij[##1]{%
    \refstepcounter{inxcnt}%
    \item[\inxlist@counterfmt{inxcnt}.]##1\hspace{.5ex}%
  }%
  \def\xisn{\item[]}%
  \def\xisnj[##1]{\item[]##1\hspace{.5ex}}%
}

\def\inxlist@close{%
  \end{list}%
  \@listdepth=\inxlist@savedlistdepth\relax
  \@itemdepth=\inxlist@saveditemdepth\relax
  \ignorespacesafterend
}

% inxlist  — alphabetic:  a. b. c.   (mirrors gb4e's xlist / xlista)
\newenvironment{inxlist}
  {\inxlist@open{\alph}{m.}}
  {\inxlist@close}

% inxlisti — roman:       i. ii. iii.  (mirrors gb4e's xlisti)
\newenvironment{inxlisti}
  {\inxlist@open{\roman}{iii.}}
  {\inxlist@close}

% inxlistn — arabic:      1. 2. 3.   (mirrors gb4e's xlistn)
\newenvironment{inxlistn}
  {\inxlist@open{\arabic}{9.}}
  {\inxlist@close}

% inxlistA — Alphabetic:  A. B. C.   (mirrors gb4e's xlistA)
\newenvironment{inxlistA}
  {\inxlist@open{\Alph}{M.}}
  {\inxlist@close}

% inxlistI — Roman:       I. II. III.  (mirrors gb4e's xlistI)
\newenvironment{inxlistI}
  {\inxlist@open{\Roman}{III.}}
  {\inxlist@close}

%% ---------------------------------------------------------------
%% inexe environment
%% ---------------------------------------------------------------

\newif\if@inexe@firstline

\newenvironment{inexe}{%
  \inexe@savedleft=\@totalleftmargin
  \advance\inexe@savedleft\leftskip
  \ifnum\value{inexecnt}<9
    \settowidth{\inexe@labelwidth}{\inexe@fmt{9}}%
  \else\ifnum\value{inexecnt}<99
    \settowidth{\inexe@labelwidth}{\inexe@fmt{99}}%
  \else
    \settowidth{\inexe@labelwidth}{\inexe@fmt{999}}%
  \fi\fi
  \@inexe@firstlinetrue
  \leavevmode\newline\noindent
  \let\ex\inexe@ex
  \let\exj\inexe@exj
  \let\sn\inexe@sn
  \let\snj\inexe@snj
  \let\exi\inexe@exi
}{%
  \vspace{-\inexe@itemsep}%
  \vspace{\belowdisplayskip}%
  \ignorespacesafterend
}

\makeatother

\endinput
 in your preamble.
%%
%% Configuration (set these BEFORE %% inexe.tex
%% A display-math-style linguistic example environment.
%% Load with %% inexe.tex
%% A display-math-style linguistic example environment.
%% Load with \input{inexe} in your preamble.
%%
%% Configuration (set these BEFORE \input{inexe}):
%%
%%   \def\inexedelimiter{parens}   % or: brackets, period
%%   \newlength{\inexelabelsep}
%%   \setlength{\inexelabelsep}{1em}
%%
%% Counter sharing with gb4e:
%%   If gb4e is loaded, inexe automatically shares gb4e's \exx counter.
%%
%% Usage:
%%   \begin{inexe}
%%     \ex{Consider sub-examples:
%%       \begin{inxlist}     % a. b. c.  (default, like xlist)
%%         \xi First sub-example.
%%         \xij[*] Bad sub-example.
%%         \xisn Unnumbered sub-example.
%%       \end{inxlist}}
%%     \ex{Roman sub-examples:
%%       \begin{inxlisti}    % i. ii. iii.  (like xlisti)
%%         \xi First.
%%         \xi Second.
%%       \end{inxlisti}}
%%   \end{inexe}
%%
%% Inxlist variants (all support \xi, \xij, \xisn, \xisnj):
%%   inxlist   a. b. c.      (mirrors xlist/xlista)
%%   inxlisti  i. ii. iii.   (mirrors xlisti)
%%   inxlistn  1. 2. 3.      (mirrors xlistn)
%%   inxlistA  A. B. C.      (mirrors xlistA)
%%   inxlistI  I. II. III.   (mirrors xlistI)

%% ---------------------------------------------------------------
%% Guard against double-loading
%% ---------------------------------------------------------------
\expandafter\ifx\csname inexe@loaded\endcsname\relax
\expandafter\def\csname inexe@loaded\endcsname{}
\else
  \expandafter\endinput
\fi

\makeatletter

%% ---------------------------------------------------------------
%% User-configurable defaults
%% (Only set if not already defined by the user before \input.)
%% ---------------------------------------------------------------

\@ifundefined{inexedelimiter}{\def\inexedelimiter{parens}}{}

\@ifundefined{inexelabelsep}{%
  \newlength{\inexelabelsep}%
  \setlength{\inexelabelsep}{1em}%
}{}

%% ---------------------------------------------------------------
%% Label formatting
%% ---------------------------------------------------------------

\def\inexe@fmt#1{%
  \ifx\inexedelimiter\@inexe@parens (#1)\else
  \ifx\inexedelimiter\@inexe@brackets [#1]\else
  \ifx\inexedelimiter\@inexe@period #1.\else
  (#1)%
  \fi\fi\fi
}
\def\@inexe@parens{parens}
\def\@inexe@brackets{brackets}
\def\@inexe@period{period}

%% ---------------------------------------------------------------
%% Counter
%% ---------------------------------------------------------------

\@definecounter{inexecnt}

\def\inexe@stepcounter{%
  \refstepcounter{inexecnt}%
  \protected@edef\@currentlabel{\arabic{inexecnt}}%
}
\def\inexe@thecounter{\arabic{inexecnt}}

\AtBeginDocument{%
  \@ifundefined{c@exx}{}{%
    \def\inexe@stepcounter{%
      \refstepcounter{exx}%
      \setcounter{inexecnt}{\value{exx}}%
      \protected@edef\@currentlabel{\arabic{exx}}%
    }%
    \def\inexe@thecounter{\arabic{exx}}%
    \def\theinexecnt{\arabic{exx}}%
  }%
}

%% ---------------------------------------------------------------
%% Core display line
%% ---------------------------------------------------------------

\newlength{\inexe@savedleft}
\newlength{\inexe@labelwidth}
\newlength{\inexe@itemsep}
\setlength{\inexe@itemsep}{3\p@ plus2\p@}

\long\def\inexe@line#1#2{%
  \if@inexe@firstline
    \@inexe@firstlinefalse
  \else
    \par
  \fi
  \noindent
  \hspace*{-\inexe@savedleft}%
  \noindent
  \hbox to \dimexpr\linewidth+\inexe@savedleft\relax{%
    \hbox to \inexe@labelwidth{%
      \hfil
      \ifx\relax#1\relax
      \else
        \inexe@fmt{#1}%
      \fi
    }%
    \hspace{\inexelabelsep}%
    \parbox[t]{\dimexpr
        \linewidth + \inexe@savedleft
        - \inexe@labelwidth
        - \inexelabelsep
      \relax}{#2}%
    \hss
  }%
  \par
  \vspace{\inexe@itemsep}%
}

%% ---------------------------------------------------------------
%% Item commands
%% ---------------------------------------------------------------

\long\def\inexe@ex#1{%
  \inexe@stepcounter
  \inexe@line{\inexe@thecounter}{#1}%
}

\long\def\inexe@exj[#1]#2{%
  \inexe@stepcounter
  \inexe@line{\inexe@thecounter}{#1\hspace{.5ex}#2}%
}

\long\def\inexe@sn#1{%
  \inexe@line{\relax}{#1}%
}

\long\def\inexe@snj[#1]#2{%
  \inexe@line{\relax}{#1\hspace{.5ex}#2}%
}

\long\def\inexe@exi#1#2{%
  \inexe@line{#1}{#2}%
}

%% ---------------------------------------------------------------
%% inxlist — inline sub-example list for use inside \ex{...}
%%
%% Lives inside the \parbox of \ex{...}. Uses \halign so \xi takes
%% unbraced content (everything to the next \xi/\xij/\xisn/\xisnj
%% or \end{inxlist}), matching gb4e's xlist feel.
%% Labels are always a. b. c. style (gb4e xlist default).
%%
%% Usage:
%%   \ex{Consider the following:
%%     \begin{inxlist}
%%       \xi First sub-example.
%%       \xi Second sub-example.
%%       \xij[*] Bad sub-example.
%%       \xisn Unnumbered sub-example.
%%     \end{inxlist}}
%%
%% Commands inside inxlist:
%%   \xi          lettered sub-example (unbraced)
%%   \xij[jdg]    with judgement mark
%%   \xisn        unnumbered sub-example
%%   \xisnj[jdg]  unnumbered with judgement
%% ---------------------------------------------------------------

\@definecounter{inxcnt}
% \p@inxcnt: prefix for \ref — produces "3b" style cross-references.
% Defined as the current inexe label (arabic) so sub-example labels
% inherit the parent example number.
\def\p@inxcnt{\theinexecnt}
% \theinexecnt: the current top-level example number, used as prefix.
% Points to exx if gb4e loaded, else inexecnt. Defined at \begin{document}
% (after gb4e has loaded) via \AtBeginDocument below — placeholder for now.
\def\theinexecnt{\arabic{inexecnt}}
% Save the standard \@item before any package (e.g. paralist) redefines it.
\let\@orignalitem\@item

\newlength{\inxlist@labelwidth}
\newlength{\inxlist@labelsep}
\setlength{\inxlist@labelsep}{.75em}
\newlength{\inxlist@itemsep}
\setlength{\inxlist@itemsep}{1.5\p@ plus\p@}

% \theinxcnt: formats the sub-example counter value (letter/roman/etc.)
% Set locally by \inxlist@open to match the current list style.
\def\theinxcnt{\alph{inxcnt}}
%% Shared setup for all inxlist variants.
%% #1 = counter format command, e.g. \alph, \roman, \arabic, \Alph, \Roman
%% #2 = sample string for label width, e.g. "m." or "iii." or "99."
\def\inxlist@open#1#2{%
  \setcounter{inxcnt}{0}%
  \def\inxlist@counterfmt{#1}%
  % Set \theinxcnt so \refstepcounter{inxcnt} records the right label.
  \def\theinxcnt{#1{inxcnt}}%
  \settowidth{\inxlist@labelwidth}{#2}%
  \edef\inxlist@savedlistdepth{\the\@listdepth}%
  \edef\inxlist@saveditemdepth{\the\@itemdepth}%
  \let\@item\@orignalitem
  \begin{list}{}{%
    \labelwidth=\inxlist@labelwidth
    \labelsep=\inxlist@labelsep
    \leftmargin=\inxlist@labelwidth
    \advance\leftmargin\inxlist@labelsep
    \itemindent\z@
    \listparindent\z@
    \topsep=\inxlist@itemsep
    \itemsep=\inxlist@itemsep
    \parsep\z@
  }%
  \def\xi{%
    \refstepcounter{inxcnt}%
    \item[\inxlist@counterfmt{inxcnt}.]%
  }%
  \def\xij[##1]{%
    \refstepcounter{inxcnt}%
    \item[\inxlist@counterfmt{inxcnt}.]##1\hspace{.5ex}%
  }%
  \def\xisn{\item[]}%
  \def\xisnj[##1]{\item[]##1\hspace{.5ex}}%
}

\def\inxlist@close{%
  \end{list}%
  \@listdepth=\inxlist@savedlistdepth\relax
  \@itemdepth=\inxlist@saveditemdepth\relax
  \ignorespacesafterend
}

% inxlist  — alphabetic:  a. b. c.   (mirrors gb4e's xlist / xlista)
\newenvironment{inxlist}
  {\inxlist@open{\alph}{m.}}
  {\inxlist@close}

% inxlisti — roman:       i. ii. iii.  (mirrors gb4e's xlisti)
\newenvironment{inxlisti}
  {\inxlist@open{\roman}{iii.}}
  {\inxlist@close}

% inxlistn — arabic:      1. 2. 3.   (mirrors gb4e's xlistn)
\newenvironment{inxlistn}
  {\inxlist@open{\arabic}{9.}}
  {\inxlist@close}

% inxlistA — Alphabetic:  A. B. C.   (mirrors gb4e's xlistA)
\newenvironment{inxlistA}
  {\inxlist@open{\Alph}{M.}}
  {\inxlist@close}

% inxlistI — Roman:       I. II. III.  (mirrors gb4e's xlistI)
\newenvironment{inxlistI}
  {\inxlist@open{\Roman}{III.}}
  {\inxlist@close}

%% ---------------------------------------------------------------
%% inexe environment
%% ---------------------------------------------------------------

\newif\if@inexe@firstline

\newenvironment{inexe}{%
  \inexe@savedleft=\@totalleftmargin
  \advance\inexe@savedleft\leftskip
  \ifnum\value{inexecnt}<9
    \settowidth{\inexe@labelwidth}{\inexe@fmt{9}}%
  \else\ifnum\value{inexecnt}<99
    \settowidth{\inexe@labelwidth}{\inexe@fmt{99}}%
  \else
    \settowidth{\inexe@labelwidth}{\inexe@fmt{999}}%
  \fi\fi
  \@inexe@firstlinetrue
  \leavevmode\newline\noindent
  \let\ex\inexe@ex
  \let\exj\inexe@exj
  \let\sn\inexe@sn
  \let\snj\inexe@snj
  \let\exi\inexe@exi
}{%
  \vspace{-\inexe@itemsep}%
  \vspace{\belowdisplayskip}%
  \ignorespacesafterend
}

\makeatother

\endinput
 in your preamble.
%%
%% Configuration (set these BEFORE %% inexe.tex
%% A display-math-style linguistic example environment.
%% Load with \input{inexe} in your preamble.
%%
%% Configuration (set these BEFORE \input{inexe}):
%%
%%   \def\inexedelimiter{parens}   % or: brackets, period
%%   \newlength{\inexelabelsep}
%%   \setlength{\inexelabelsep}{1em}
%%
%% Counter sharing with gb4e:
%%   If gb4e is loaded, inexe automatically shares gb4e's \exx counter.
%%
%% Usage:
%%   \begin{inexe}
%%     \ex{Consider sub-examples:
%%       \begin{inxlist}     % a. b. c.  (default, like xlist)
%%         \xi First sub-example.
%%         \xij[*] Bad sub-example.
%%         \xisn Unnumbered sub-example.
%%       \end{inxlist}}
%%     \ex{Roman sub-examples:
%%       \begin{inxlisti}    % i. ii. iii.  (like xlisti)
%%         \xi First.
%%         \xi Second.
%%       \end{inxlisti}}
%%   \end{inexe}
%%
%% Inxlist variants (all support \xi, \xij, \xisn, \xisnj):
%%   inxlist   a. b. c.      (mirrors xlist/xlista)
%%   inxlisti  i. ii. iii.   (mirrors xlisti)
%%   inxlistn  1. 2. 3.      (mirrors xlistn)
%%   inxlistA  A. B. C.      (mirrors xlistA)
%%   inxlistI  I. II. III.   (mirrors xlistI)

%% ---------------------------------------------------------------
%% Guard against double-loading
%% ---------------------------------------------------------------
\expandafter\ifx\csname inexe@loaded\endcsname\relax
\expandafter\def\csname inexe@loaded\endcsname{}
\else
  \expandafter\endinput
\fi

\makeatletter

%% ---------------------------------------------------------------
%% User-configurable defaults
%% (Only set if not already defined by the user before \input.)
%% ---------------------------------------------------------------

\@ifundefined{inexedelimiter}{\def\inexedelimiter{parens}}{}

\@ifundefined{inexelabelsep}{%
  \newlength{\inexelabelsep}%
  \setlength{\inexelabelsep}{1em}%
}{}

%% ---------------------------------------------------------------
%% Label formatting
%% ---------------------------------------------------------------

\def\inexe@fmt#1{%
  \ifx\inexedelimiter\@inexe@parens (#1)\else
  \ifx\inexedelimiter\@inexe@brackets [#1]\else
  \ifx\inexedelimiter\@inexe@period #1.\else
  (#1)%
  \fi\fi\fi
}
\def\@inexe@parens{parens}
\def\@inexe@brackets{brackets}
\def\@inexe@period{period}

%% ---------------------------------------------------------------
%% Counter
%% ---------------------------------------------------------------

\@definecounter{inexecnt}

\def\inexe@stepcounter{%
  \refstepcounter{inexecnt}%
  \protected@edef\@currentlabel{\arabic{inexecnt}}%
}
\def\inexe@thecounter{\arabic{inexecnt}}

\AtBeginDocument{%
  \@ifundefined{c@exx}{}{%
    \def\inexe@stepcounter{%
      \refstepcounter{exx}%
      \setcounter{inexecnt}{\value{exx}}%
      \protected@edef\@currentlabel{\arabic{exx}}%
    }%
    \def\inexe@thecounter{\arabic{exx}}%
    \def\theinexecnt{\arabic{exx}}%
  }%
}

%% ---------------------------------------------------------------
%% Core display line
%% ---------------------------------------------------------------

\newlength{\inexe@savedleft}
\newlength{\inexe@labelwidth}
\newlength{\inexe@itemsep}
\setlength{\inexe@itemsep}{3\p@ plus2\p@}

\long\def\inexe@line#1#2{%
  \if@inexe@firstline
    \@inexe@firstlinefalse
  \else
    \par
  \fi
  \noindent
  \hspace*{-\inexe@savedleft}%
  \noindent
  \hbox to \dimexpr\linewidth+\inexe@savedleft\relax{%
    \hbox to \inexe@labelwidth{%
      \hfil
      \ifx\relax#1\relax
      \else
        \inexe@fmt{#1}%
      \fi
    }%
    \hspace{\inexelabelsep}%
    \parbox[t]{\dimexpr
        \linewidth + \inexe@savedleft
        - \inexe@labelwidth
        - \inexelabelsep
      \relax}{#2}%
    \hss
  }%
  \par
  \vspace{\inexe@itemsep}%
}

%% ---------------------------------------------------------------
%% Item commands
%% ---------------------------------------------------------------

\long\def\inexe@ex#1{%
  \inexe@stepcounter
  \inexe@line{\inexe@thecounter}{#1}%
}

\long\def\inexe@exj[#1]#2{%
  \inexe@stepcounter
  \inexe@line{\inexe@thecounter}{#1\hspace{.5ex}#2}%
}

\long\def\inexe@sn#1{%
  \inexe@line{\relax}{#1}%
}

\long\def\inexe@snj[#1]#2{%
  \inexe@line{\relax}{#1\hspace{.5ex}#2}%
}

\long\def\inexe@exi#1#2{%
  \inexe@line{#1}{#2}%
}

%% ---------------------------------------------------------------
%% inxlist — inline sub-example list for use inside \ex{...}
%%
%% Lives inside the \parbox of \ex{...}. Uses \halign so \xi takes
%% unbraced content (everything to the next \xi/\xij/\xisn/\xisnj
%% or \end{inxlist}), matching gb4e's xlist feel.
%% Labels are always a. b. c. style (gb4e xlist default).
%%
%% Usage:
%%   \ex{Consider the following:
%%     \begin{inxlist}
%%       \xi First sub-example.
%%       \xi Second sub-example.
%%       \xij[*] Bad sub-example.
%%       \xisn Unnumbered sub-example.
%%     \end{inxlist}}
%%
%% Commands inside inxlist:
%%   \xi          lettered sub-example (unbraced)
%%   \xij[jdg]    with judgement mark
%%   \xisn        unnumbered sub-example
%%   \xisnj[jdg]  unnumbered with judgement
%% ---------------------------------------------------------------

\@definecounter{inxcnt}
% \p@inxcnt: prefix for \ref — produces "3b" style cross-references.
% Defined as the current inexe label (arabic) so sub-example labels
% inherit the parent example number.
\def\p@inxcnt{\theinexecnt}
% \theinexecnt: the current top-level example number, used as prefix.
% Points to exx if gb4e loaded, else inexecnt. Defined at \begin{document}
% (after gb4e has loaded) via \AtBeginDocument below — placeholder for now.
\def\theinexecnt{\arabic{inexecnt}}
% Save the standard \@item before any package (e.g. paralist) redefines it.
\let\@orignalitem\@item

\newlength{\inxlist@labelwidth}
\newlength{\inxlist@labelsep}
\setlength{\inxlist@labelsep}{.75em}
\newlength{\inxlist@itemsep}
\setlength{\inxlist@itemsep}{1.5\p@ plus\p@}

% \theinxcnt: formats the sub-example counter value (letter/roman/etc.)
% Set locally by \inxlist@open to match the current list style.
\def\theinxcnt{\alph{inxcnt}}
%% Shared setup for all inxlist variants.
%% #1 = counter format command, e.g. \alph, \roman, \arabic, \Alph, \Roman
%% #2 = sample string for label width, e.g. "m." or "iii." or "99."
\def\inxlist@open#1#2{%
  \setcounter{inxcnt}{0}%
  \def\inxlist@counterfmt{#1}%
  % Set \theinxcnt so \refstepcounter{inxcnt} records the right label.
  \def\theinxcnt{#1{inxcnt}}%
  \settowidth{\inxlist@labelwidth}{#2}%
  \edef\inxlist@savedlistdepth{\the\@listdepth}%
  \edef\inxlist@saveditemdepth{\the\@itemdepth}%
  \let\@item\@orignalitem
  \begin{list}{}{%
    \labelwidth=\inxlist@labelwidth
    \labelsep=\inxlist@labelsep
    \leftmargin=\inxlist@labelwidth
    \advance\leftmargin\inxlist@labelsep
    \itemindent\z@
    \listparindent\z@
    \topsep=\inxlist@itemsep
    \itemsep=\inxlist@itemsep
    \parsep\z@
  }%
  \def\xi{%
    \refstepcounter{inxcnt}%
    \item[\inxlist@counterfmt{inxcnt}.]%
  }%
  \def\xij[##1]{%
    \refstepcounter{inxcnt}%
    \item[\inxlist@counterfmt{inxcnt}.]##1\hspace{.5ex}%
  }%
  \def\xisn{\item[]}%
  \def\xisnj[##1]{\item[]##1\hspace{.5ex}}%
}

\def\inxlist@close{%
  \end{list}%
  \@listdepth=\inxlist@savedlistdepth\relax
  \@itemdepth=\inxlist@saveditemdepth\relax
  \ignorespacesafterend
}

% inxlist  — alphabetic:  a. b. c.   (mirrors gb4e's xlist / xlista)
\newenvironment{inxlist}
  {\inxlist@open{\alph}{m.}}
  {\inxlist@close}

% inxlisti — roman:       i. ii. iii.  (mirrors gb4e's xlisti)
\newenvironment{inxlisti}
  {\inxlist@open{\roman}{iii.}}
  {\inxlist@close}

% inxlistn — arabic:      1. 2. 3.   (mirrors gb4e's xlistn)
\newenvironment{inxlistn}
  {\inxlist@open{\arabic}{9.}}
  {\inxlist@close}

% inxlistA — Alphabetic:  A. B. C.   (mirrors gb4e's xlistA)
\newenvironment{inxlistA}
  {\inxlist@open{\Alph}{M.}}
  {\inxlist@close}

% inxlistI — Roman:       I. II. III.  (mirrors gb4e's xlistI)
\newenvironment{inxlistI}
  {\inxlist@open{\Roman}{III.}}
  {\inxlist@close}

%% ---------------------------------------------------------------
%% inexe environment
%% ---------------------------------------------------------------

\newif\if@inexe@firstline

\newenvironment{inexe}{%
  \inexe@savedleft=\@totalleftmargin
  \advance\inexe@savedleft\leftskip
  \ifnum\value{inexecnt}<9
    \settowidth{\inexe@labelwidth}{\inexe@fmt{9}}%
  \else\ifnum\value{inexecnt}<99
    \settowidth{\inexe@labelwidth}{\inexe@fmt{99}}%
  \else
    \settowidth{\inexe@labelwidth}{\inexe@fmt{999}}%
  \fi\fi
  \@inexe@firstlinetrue
  \leavevmode\newline\noindent
  \let\ex\inexe@ex
  \let\exj\inexe@exj
  \let\sn\inexe@sn
  \let\snj\inexe@snj
  \let\exi\inexe@exi
}{%
  \vspace{-\inexe@itemsep}%
  \vspace{\belowdisplayskip}%
  \ignorespacesafterend
}

\makeatother

\endinput
):
%%
%%   \def\inexedelimiter{parens}   % or: brackets, period
%%   \newlength{\inexelabelsep}
%%   \setlength{\inexelabelsep}{1em}
%%
%% Counter sharing with gb4e:
%%   If gb4e is loaded, inexe automatically shares gb4e's \exx counter.
%%
%% Usage:
%%   \begin{inexe}
%%     \ex{Consider sub-examples:
%%       \begin{inxlist}     % a. b. c.  (default, like xlist)
%%         \xi First sub-example.
%%         \xij[*] Bad sub-example.
%%         \xisn Unnumbered sub-example.
%%       \end{inxlist}}
%%     \ex{Roman sub-examples:
%%       \begin{inxlisti}    % i. ii. iii.  (like xlisti)
%%         \xi First.
%%         \xi Second.
%%       \end{inxlisti}}
%%   \end{inexe}
%%
%% Inxlist variants (all support \xi, \xij, \xisn, \xisnj):
%%   inxlist   a. b. c.      (mirrors xlist/xlista)
%%   inxlisti  i. ii. iii.   (mirrors xlisti)
%%   inxlistn  1. 2. 3.      (mirrors xlistn)
%%   inxlistA  A. B. C.      (mirrors xlistA)
%%   inxlistI  I. II. III.   (mirrors xlistI)

%% ---------------------------------------------------------------
%% Guard against double-loading
%% ---------------------------------------------------------------
\expandafter\ifx\csname inexe@loaded\endcsname\relax
\expandafter\def\csname inexe@loaded\endcsname{}
\else
  \expandafter\endinput
\fi

\makeatletter

%% ---------------------------------------------------------------
%% User-configurable defaults
%% (Only set if not already defined by the user before \input.)
%% ---------------------------------------------------------------

\@ifundefined{inexedelimiter}{\def\inexedelimiter{parens}}{}

\@ifundefined{inexelabelsep}{%
  \newlength{\inexelabelsep}%
  \setlength{\inexelabelsep}{1em}%
}{}

%% ---------------------------------------------------------------
%% Label formatting
%% ---------------------------------------------------------------

\def\inexe@fmt#1{%
  \ifx\inexedelimiter\@inexe@parens (#1)\else
  \ifx\inexedelimiter\@inexe@brackets [#1]\else
  \ifx\inexedelimiter\@inexe@period #1.\else
  (#1)%
  \fi\fi\fi
}
\def\@inexe@parens{parens}
\def\@inexe@brackets{brackets}
\def\@inexe@period{period}

%% ---------------------------------------------------------------
%% Counter
%% ---------------------------------------------------------------

\@definecounter{inexecnt}

\def\inexe@stepcounter{%
  \refstepcounter{inexecnt}%
  \protected@edef\@currentlabel{\arabic{inexecnt}}%
}
\def\inexe@thecounter{\arabic{inexecnt}}

\AtBeginDocument{%
  \@ifundefined{c@exx}{}{%
    \def\inexe@stepcounter{%
      \refstepcounter{exx}%
      \setcounter{inexecnt}{\value{exx}}%
      \protected@edef\@currentlabel{\arabic{exx}}%
    }%
    \def\inexe@thecounter{\arabic{exx}}%
    \def\theinexecnt{\arabic{exx}}%
  }%
}

%% ---------------------------------------------------------------
%% Core display line
%% ---------------------------------------------------------------

\newlength{\inexe@savedleft}
\newlength{\inexe@labelwidth}
\newlength{\inexe@itemsep}
\setlength{\inexe@itemsep}{3\p@ plus2\p@}

\long\def\inexe@line#1#2{%
  \if@inexe@firstline
    \@inexe@firstlinefalse
  \else
    \par
  \fi
  \noindent
  \hspace*{-\inexe@savedleft}%
  \noindent
  \hbox to \dimexpr\linewidth+\inexe@savedleft\relax{%
    \hbox to \inexe@labelwidth{%
      \hfil
      \ifx\relax#1\relax
      \else
        \inexe@fmt{#1}%
      \fi
    }%
    \hspace{\inexelabelsep}%
    \parbox[t]{\dimexpr
        \linewidth + \inexe@savedleft
        - \inexe@labelwidth
        - \inexelabelsep
      \relax}{#2}%
    \hss
  }%
  \par
  \vspace{\inexe@itemsep}%
}

%% ---------------------------------------------------------------
%% Item commands
%% ---------------------------------------------------------------

\long\def\inexe@ex#1{%
  \inexe@stepcounter
  \inexe@line{\inexe@thecounter}{#1}%
}

\long\def\inexe@exj[#1]#2{%
  \inexe@stepcounter
  \inexe@line{\inexe@thecounter}{#1\hspace{.5ex}#2}%
}

\long\def\inexe@sn#1{%
  \inexe@line{\relax}{#1}%
}

\long\def\inexe@snj[#1]#2{%
  \inexe@line{\relax}{#1\hspace{.5ex}#2}%
}

\long\def\inexe@exi#1#2{%
  \inexe@line{#1}{#2}%
}

%% ---------------------------------------------------------------
%% inxlist — inline sub-example list for use inside \ex{...}
%%
%% Lives inside the \parbox of \ex{...}. Uses \halign so \xi takes
%% unbraced content (everything to the next \xi/\xij/\xisn/\xisnj
%% or \end{inxlist}), matching gb4e's xlist feel.
%% Labels are always a. b. c. style (gb4e xlist default).
%%
%% Usage:
%%   \ex{Consider the following:
%%     \begin{inxlist}
%%       \xi First sub-example.
%%       \xi Second sub-example.
%%       \xij[*] Bad sub-example.
%%       \xisn Unnumbered sub-example.
%%     \end{inxlist}}
%%
%% Commands inside inxlist:
%%   \xi          lettered sub-example (unbraced)
%%   \xij[jdg]    with judgement mark
%%   \xisn        unnumbered sub-example
%%   \xisnj[jdg]  unnumbered with judgement
%% ---------------------------------------------------------------

\@definecounter{inxcnt}
% \p@inxcnt: prefix for \ref — produces "3b" style cross-references.
% Defined as the current inexe label (arabic) so sub-example labels
% inherit the parent example number.
\def\p@inxcnt{\theinexecnt}
% \theinexecnt: the current top-level example number, used as prefix.
% Points to exx if gb4e loaded, else inexecnt. Defined at \begin{document}
% (after gb4e has loaded) via \AtBeginDocument below — placeholder for now.
\def\theinexecnt{\arabic{inexecnt}}
% Save the standard \@item before any package (e.g. paralist) redefines it.
\let\@orignalitem\@item

\newlength{\inxlist@labelwidth}
\newlength{\inxlist@labelsep}
\setlength{\inxlist@labelsep}{.75em}
\newlength{\inxlist@itemsep}
\setlength{\inxlist@itemsep}{1.5\p@ plus\p@}

% \theinxcnt: formats the sub-example counter value (letter/roman/etc.)
% Set locally by \inxlist@open to match the current list style.
\def\theinxcnt{\alph{inxcnt}}
%% Shared setup for all inxlist variants.
%% #1 = counter format command, e.g. \alph, \roman, \arabic, \Alph, \Roman
%% #2 = sample string for label width, e.g. "m." or "iii." or "99."
\def\inxlist@open#1#2{%
  \setcounter{inxcnt}{0}%
  \def\inxlist@counterfmt{#1}%
  % Set \theinxcnt so \refstepcounter{inxcnt} records the right label.
  \def\theinxcnt{#1{inxcnt}}%
  \settowidth{\inxlist@labelwidth}{#2}%
  \edef\inxlist@savedlistdepth{\the\@listdepth}%
  \edef\inxlist@saveditemdepth{\the\@itemdepth}%
  \let\@item\@orignalitem
  \begin{list}{}{%
    \labelwidth=\inxlist@labelwidth
    \labelsep=\inxlist@labelsep
    \leftmargin=\inxlist@labelwidth
    \advance\leftmargin\inxlist@labelsep
    \itemindent\z@
    \listparindent\z@
    \topsep=\inxlist@itemsep
    \itemsep=\inxlist@itemsep
    \parsep\z@
  }%
  \def\xi{%
    \refstepcounter{inxcnt}%
    \item[\inxlist@counterfmt{inxcnt}.]%
  }%
  \def\xij[##1]{%
    \refstepcounter{inxcnt}%
    \item[\inxlist@counterfmt{inxcnt}.]##1\hspace{.5ex}%
  }%
  \def\xisn{\item[]}%
  \def\xisnj[##1]{\item[]##1\hspace{.5ex}}%
}

\def\inxlist@close{%
  \end{list}%
  \@listdepth=\inxlist@savedlistdepth\relax
  \@itemdepth=\inxlist@saveditemdepth\relax
  \ignorespacesafterend
}

% inxlist  — alphabetic:  a. b. c.   (mirrors gb4e's xlist / xlista)
\newenvironment{inxlist}
  {\inxlist@open{\alph}{m.}}
  {\inxlist@close}

% inxlisti — roman:       i. ii. iii.  (mirrors gb4e's xlisti)
\newenvironment{inxlisti}
  {\inxlist@open{\roman}{iii.}}
  {\inxlist@close}

% inxlistn — arabic:      1. 2. 3.   (mirrors gb4e's xlistn)
\newenvironment{inxlistn}
  {\inxlist@open{\arabic}{9.}}
  {\inxlist@close}

% inxlistA — Alphabetic:  A. B. C.   (mirrors gb4e's xlistA)
\newenvironment{inxlistA}
  {\inxlist@open{\Alph}{M.}}
  {\inxlist@close}

% inxlistI — Roman:       I. II. III.  (mirrors gb4e's xlistI)
\newenvironment{inxlistI}
  {\inxlist@open{\Roman}{III.}}
  {\inxlist@close}

%% ---------------------------------------------------------------
%% inexe environment
%% ---------------------------------------------------------------

\newif\if@inexe@firstline

\newenvironment{inexe}{%
  \inexe@savedleft=\@totalleftmargin
  \advance\inexe@savedleft\leftskip
  \ifnum\value{inexecnt}<9
    \settowidth{\inexe@labelwidth}{\inexe@fmt{9}}%
  \else\ifnum\value{inexecnt}<99
    \settowidth{\inexe@labelwidth}{\inexe@fmt{99}}%
  \else
    \settowidth{\inexe@labelwidth}{\inexe@fmt{999}}%
  \fi\fi
  \@inexe@firstlinetrue
  \leavevmode\newline\noindent
  \let\ex\inexe@ex
  \let\exj\inexe@exj
  \let\sn\inexe@sn
  \let\snj\inexe@snj
  \let\exi\inexe@exi
}{%
  \vspace{-\inexe@itemsep}%
  \vspace{\belowdisplayskip}%
  \ignorespacesafterend
}

\makeatother

\endinput
):
%%
%%   \def\inexedelimiter{parens}   % or: brackets, period
%%   \newlength{\inexelabelsep}
%%   \setlength{\inexelabelsep}{1em}
%%
%% Counter sharing with gb4e:
%%   If gb4e is loaded, inexe automatically shares gb4e's \exx counter.
%%
%% Usage:
%%   \begin{inexe}
%%     \ex{Consider sub-examples:
%%       \begin{inxlist}     % a. b. c.  (default, like xlist)
%%         \xi First sub-example.
%%         \xij[*] Bad sub-example.
%%         \xisn Unnumbered sub-example.
%%       \end{inxlist}}
%%     \ex{Roman sub-examples:
%%       \begin{inxlisti}    % i. ii. iii.  (like xlisti)
%%         \xi First.
%%         \xi Second.
%%       \end{inxlisti}}
%%   \end{inexe}
%%
%% Inxlist variants (all support \xi, \xij, \xisn, \xisnj):
%%   inxlist   a. b. c.      (mirrors xlist/xlista)
%%   inxlisti  i. ii. iii.   (mirrors xlisti)
%%   inxlistn  1. 2. 3.      (mirrors xlistn)
%%   inxlistA  A. B. C.      (mirrors xlistA)
%%   inxlistI  I. II. III.   (mirrors xlistI)

%% ---------------------------------------------------------------
%% Guard against double-loading
%% ---------------------------------------------------------------
\expandafter\ifx\csname inexe@loaded\endcsname\relax
\expandafter\def\csname inexe@loaded\endcsname{}
\else
  \expandafter\endinput
\fi

\makeatletter

%% ---------------------------------------------------------------
%% User-configurable defaults
%% (Only set if not already defined by the user before \input.)
%% ---------------------------------------------------------------

\@ifundefined{inexedelimiter}{\def\inexedelimiter{parens}}{}

\@ifundefined{inexelabelsep}{%
  \newlength{\inexelabelsep}%
  \setlength{\inexelabelsep}{1em}%
}{}

%% ---------------------------------------------------------------
%% Label formatting
%% ---------------------------------------------------------------

\def\inexe@fmt#1{%
  \ifx\inexedelimiter\@inexe@parens (#1)\else
  \ifx\inexedelimiter\@inexe@brackets [#1]\else
  \ifx\inexedelimiter\@inexe@period #1.\else
  (#1)%
  \fi\fi\fi
}
\def\@inexe@parens{parens}
\def\@inexe@brackets{brackets}
\def\@inexe@period{period}

%% ---------------------------------------------------------------
%% Counter
%% ---------------------------------------------------------------

\@definecounter{inexecnt}

\def\inexe@stepcounter{%
  \refstepcounter{inexecnt}%
  \protected@edef\@currentlabel{\arabic{inexecnt}}%
}
\def\inexe@thecounter{\arabic{inexecnt}}

\AtBeginDocument{%
  \@ifundefined{c@exx}{}{%
    \def\inexe@stepcounter{%
      \refstepcounter{exx}%
      \setcounter{inexecnt}{\value{exx}}%
      \protected@edef\@currentlabel{\arabic{exx}}%
    }%
    \def\inexe@thecounter{\arabic{exx}}%
    \def\theinexecnt{\arabic{exx}}%
  }%
}

%% ---------------------------------------------------------------
%% Core display line
%% ---------------------------------------------------------------

\newlength{\inexe@savedleft}
\newlength{\inexe@labelwidth}
\newlength{\inexe@itemsep}
\setlength{\inexe@itemsep}{3\p@ plus2\p@}

\long\def\inexe@line#1#2{%
  \if@inexe@firstline
    \@inexe@firstlinefalse
  \else
    \par
  \fi
  \noindent
  \hspace*{-\inexe@savedleft}%
  \noindent
  \hbox to \dimexpr\linewidth+\inexe@savedleft\relax{%
    \hbox to \inexe@labelwidth{%
      \hfil
      \ifx\relax#1\relax
      \else
        \inexe@fmt{#1}%
      \fi
    }%
    \hspace{\inexelabelsep}%
    \parbox[t]{\dimexpr
        \linewidth + \inexe@savedleft
        - \inexe@labelwidth
        - \inexelabelsep
      \relax}{#2}%
    \hss
  }%
  \par
  \vspace{\inexe@itemsep}%
}

%% ---------------------------------------------------------------
%% Item commands
%% ---------------------------------------------------------------

\long\def\inexe@ex#1{%
  \inexe@stepcounter
  \inexe@line{\inexe@thecounter}{#1}%
}

\long\def\inexe@exj[#1]#2{%
  \inexe@stepcounter
  \inexe@line{\inexe@thecounter}{#1\hspace{.5ex}#2}%
}

\long\def\inexe@sn#1{%
  \inexe@line{\relax}{#1}%
}

\long\def\inexe@snj[#1]#2{%
  \inexe@line{\relax}{#1\hspace{.5ex}#2}%
}

\long\def\inexe@exi#1#2{%
  \inexe@line{#1}{#2}%
}

%% ---------------------------------------------------------------
%% inxlist — inline sub-example list for use inside \ex{...}
%%
%% Lives inside the \parbox of \ex{...}. Uses \halign so \xi takes
%% unbraced content (everything to the next \xi/\xij/\xisn/\xisnj
%% or \end{inxlist}), matching gb4e's xlist feel.
%% Labels are always a. b. c. style (gb4e xlist default).
%%
%% Usage:
%%   \ex{Consider the following:
%%     \begin{inxlist}
%%       \xi First sub-example.
%%       \xi Second sub-example.
%%       \xij[*] Bad sub-example.
%%       \xisn Unnumbered sub-example.
%%     \end{inxlist}}
%%
%% Commands inside inxlist:
%%   \xi          lettered sub-example (unbraced)
%%   \xij[jdg]    with judgement mark
%%   \xisn        unnumbered sub-example
%%   \xisnj[jdg]  unnumbered with judgement
%% ---------------------------------------------------------------

\@definecounter{inxcnt}
% \p@inxcnt: prefix for \ref — produces "3b" style cross-references.
% Defined as the current inexe label (arabic) so sub-example labels
% inherit the parent example number.
\def\p@inxcnt{\theinexecnt}
% \theinexecnt: the current top-level example number, used as prefix.
% Points to exx if gb4e loaded, else inexecnt. Defined at \begin{document}
% (after gb4e has loaded) via \AtBeginDocument below — placeholder for now.
\def\theinexecnt{\arabic{inexecnt}}
% Save the standard \@item before any package (e.g. paralist) redefines it.
\let\@orignalitem\@item

\newlength{\inxlist@labelwidth}
\newlength{\inxlist@labelsep}
\setlength{\inxlist@labelsep}{.75em}
\newlength{\inxlist@itemsep}
\setlength{\inxlist@itemsep}{1.5\p@ plus\p@}

% \theinxcnt: formats the sub-example counter value (letter/roman/etc.)
% Set locally by \inxlist@open to match the current list style.
\def\theinxcnt{\alph{inxcnt}}
%% Shared setup for all inxlist variants.
%% #1 = counter format command, e.g. \alph, \roman, \arabic, \Alph, \Roman
%% #2 = sample string for label width, e.g. "m." or "iii." or "99."
\def\inxlist@open#1#2{%
  \setcounter{inxcnt}{0}%
  \def\inxlist@counterfmt{#1}%
  % Set \theinxcnt so \refstepcounter{inxcnt} records the right label.
  \def\theinxcnt{#1{inxcnt}}%
  \settowidth{\inxlist@labelwidth}{#2}%
  \edef\inxlist@savedlistdepth{\the\@listdepth}%
  \edef\inxlist@saveditemdepth{\the\@itemdepth}%
  \let\@item\@orignalitem
  \begin{list}{}{%
    \labelwidth=\inxlist@labelwidth
    \labelsep=\inxlist@labelsep
    \leftmargin=\inxlist@labelwidth
    \advance\leftmargin\inxlist@labelsep
    \itemindent\z@
    \listparindent\z@
    \topsep=\inxlist@itemsep
    \itemsep=\inxlist@itemsep
    \parsep\z@
  }%
  \def\xi{%
    \refstepcounter{inxcnt}%
    \item[\inxlist@counterfmt{inxcnt}.]%
  }%
  \def\xij[##1]{%
    \refstepcounter{inxcnt}%
    \item[\inxlist@counterfmt{inxcnt}.]##1\hspace{.5ex}%
  }%
  \def\xisn{\item[]}%
  \def\xisnj[##1]{\item[]##1\hspace{.5ex}}%
}

\def\inxlist@close{%
  \end{list}%
  \@listdepth=\inxlist@savedlistdepth\relax
  \@itemdepth=\inxlist@saveditemdepth\relax
  \ignorespacesafterend
}

% inxlist  — alphabetic:  a. b. c.   (mirrors gb4e's xlist / xlista)
\newenvironment{inxlist}
  {\inxlist@open{\alph}{m.}}
  {\inxlist@close}

% inxlisti — roman:       i. ii. iii.  (mirrors gb4e's xlisti)
\newenvironment{inxlisti}
  {\inxlist@open{\roman}{iii.}}
  {\inxlist@close}

% inxlistn — arabic:      1. 2. 3.   (mirrors gb4e's xlistn)
\newenvironment{inxlistn}
  {\inxlist@open{\arabic}{9.}}
  {\inxlist@close}

% inxlistA — Alphabetic:  A. B. C.   (mirrors gb4e's xlistA)
\newenvironment{inxlistA}
  {\inxlist@open{\Alph}{M.}}
  {\inxlist@close}

% inxlistI — Roman:       I. II. III.  (mirrors gb4e's xlistI)
\newenvironment{inxlistI}
  {\inxlist@open{\Roman}{III.}}
  {\inxlist@close}

%% ---------------------------------------------------------------
%% inexe environment
%% ---------------------------------------------------------------

\newif\if@inexe@firstline

\newenvironment{inexe}{%
  \inexe@savedleft=\@totalleftmargin
  \advance\inexe@savedleft\leftskip
  \ifnum\value{inexecnt}<9
    \settowidth{\inexe@labelwidth}{\inexe@fmt{9}}%
  \else\ifnum\value{inexecnt}<99
    \settowidth{\inexe@labelwidth}{\inexe@fmt{99}}%
  \else
    \settowidth{\inexe@labelwidth}{\inexe@fmt{999}}%
  \fi\fi
  \@inexe@firstlinetrue
  \leavevmode\newline\noindent
  \let\ex\inexe@ex
  \let\exj\inexe@exj
  \let\sn\inexe@sn
  \let\snj\inexe@snj
  \let\exi\inexe@exi
}{%
  \vspace{-\inexe@itemsep}%
  \vspace{\belowdisplayskip}%
  \ignorespacesafterend
}

\makeatother

\endinput
):
%%
%%   \def\inexedelimiter{parens}   % or: brackets, period
%%   \newlength{\inexelabelsep}
%%   \setlength{\inexelabelsep}{1em}
%%
%% Counter sharing with gb4e:
%%   If gb4e is loaded, inexe automatically shares gb4e's \exx counter.
%%
%% Usage:
%%   \begin{inexe}
%%     \ex{Consider sub-examples:
%%       \begin{inxlist}     % a. b. c.  (default, like xlist)
%%         \xi First sub-example.
%%         \xij[*] Bad sub-example.
%%         \xisn Unnumbered sub-example.
%%       \end{inxlist}}
%%     \ex{Roman sub-examples:
%%       \begin{inxlisti}    % i. ii. iii.  (like xlisti)
%%         \xi First.
%%         \xi Second.
%%       \end{inxlisti}}
%%   \end{inexe}
%%
%% Inxlist variants (all support \xi, \xij, \xisn, \xisnj):
%%   inxlist   a. b. c.      (mirrors xlist/xlista)
%%   inxlisti  i. ii. iii.   (mirrors xlisti)
%%   inxlistn  1. 2. 3.      (mirrors xlistn)
%%   inxlistA  A. B. C.      (mirrors xlistA)
%%   inxlistI  I. II. III.   (mirrors xlistI)

%% ---------------------------------------------------------------
%% Guard against double-loading
%% ---------------------------------------------------------------
\expandafter\ifx\csname inexe@loaded\endcsname\relax
\expandafter\def\csname inexe@loaded\endcsname{}
\else
  \expandafter\endinput
\fi

\makeatletter

%% ---------------------------------------------------------------
%% User-configurable defaults
%% (Only set if not already defined by the user before \input.)
%% ---------------------------------------------------------------

\@ifundefined{inexedelimiter}{\def\inexedelimiter{parens}}{}

\@ifundefined{inexelabelsep}{%
  \newlength{\inexelabelsep}%
  \setlength{\inexelabelsep}{1em}%
}{}

%% ---------------------------------------------------------------
%% Label formatting
%% ---------------------------------------------------------------

\def\inexe@fmt#1{%
  \ifx\inexedelimiter\@inexe@parens (#1)\else
  \ifx\inexedelimiter\@inexe@brackets [#1]\else
  \ifx\inexedelimiter\@inexe@period #1.\else
  (#1)%
  \fi\fi\fi
}
\def\@inexe@parens{parens}
\def\@inexe@brackets{brackets}
\def\@inexe@period{period}

%% ---------------------------------------------------------------
%% Counter
%% ---------------------------------------------------------------

\@definecounter{inexecnt}

\def\inexe@stepcounter{%
  \refstepcounter{inexecnt}%
  \protected@edef\@currentlabel{\arabic{inexecnt}}%
}
\def\inexe@thecounter{\arabic{inexecnt}}

\AtBeginDocument{%
  \@ifundefined{c@exx}{}{%
    \def\inexe@stepcounter{%
      \refstepcounter{exx}%
      \setcounter{inexecnt}{\value{exx}}%
      \protected@edef\@currentlabel{\arabic{exx}}%
    }%
    \def\inexe@thecounter{\arabic{exx}}%
    \def\theinexecnt{\arabic{exx}}%
  }%
}

%% ---------------------------------------------------------------
%% Core display line
%% ---------------------------------------------------------------

\newlength{\inexe@savedleft}
\newlength{\inexe@labelwidth}
\newlength{\inexe@itemsep}
\setlength{\inexe@itemsep}{3\p@ plus2\p@}

\long\def\inexe@line#1#2{%
  \if@inexe@firstline
    \@inexe@firstlinefalse
  \else
    \par
  \fi
  \noindent
  \hspace*{-\inexe@savedleft}%
  \noindent
  \hbox to \dimexpr\linewidth+\inexe@savedleft\relax{%
    \hbox to \inexe@labelwidth{%
      \hfil
      \ifx\relax#1\relax
      \else
        \inexe@fmt{#1}%
      \fi
    }%
    \hspace{\inexelabelsep}%
    \parbox[t]{\dimexpr
        \linewidth + \inexe@savedleft
        - \inexe@labelwidth
        - \inexelabelsep
      \relax}{#2}%
    \hss
  }%
  \par
  \vspace{\inexe@itemsep}%
}

%% ---------------------------------------------------------------
%% Item commands
%% ---------------------------------------------------------------

\long\def\inexe@ex#1{%
  \inexe@stepcounter
  \inexe@line{\inexe@thecounter}{#1}%
}

\long\def\inexe@exj[#1]#2{%
  \inexe@stepcounter
  \inexe@line{\inexe@thecounter}{#1\hspace{.5ex}#2}%
}

\long\def\inexe@sn#1{%
  \inexe@line{\relax}{#1}%
}

\long\def\inexe@snj[#1]#2{%
  \inexe@line{\relax}{#1\hspace{.5ex}#2}%
}

\long\def\inexe@exi#1#2{%
  \inexe@line{#1}{#2}%
}

%% ---------------------------------------------------------------
%% inxlist — inline sub-example list for use inside \ex{...}
%%
%% Lives inside the \parbox of \ex{...}. Uses \halign so \xi takes
%% unbraced content (everything to the next \xi/\xij/\xisn/\xisnj
%% or \end{inxlist}), matching gb4e's xlist feel.
%% Labels are always a. b. c. style (gb4e xlist default).
%%
%% Usage:
%%   \ex{Consider the following:
%%     \begin{inxlist}
%%       \xi First sub-example.
%%       \xi Second sub-example.
%%       \xij[*] Bad sub-example.
%%       \xisn Unnumbered sub-example.
%%     \end{inxlist}}
%%
%% Commands inside inxlist:
%%   \xi          lettered sub-example (unbraced)
%%   \xij[jdg]    with judgement mark
%%   \xisn        unnumbered sub-example
%%   \xisnj[jdg]  unnumbered with judgement
%% ---------------------------------------------------------------

\@definecounter{inxcnt}
% \p@inxcnt: prefix for \ref — produces "3b" style cross-references.
% Defined as the current inexe label (arabic) so sub-example labels
% inherit the parent example number.
\def\p@inxcnt{\theinexecnt}
% \theinexecnt: the current top-level example number, used as prefix.
% Points to exx if gb4e loaded, else inexecnt. Defined at \begin{document}
% (after gb4e has loaded) via \AtBeginDocument below — placeholder for now.
\def\theinexecnt{\arabic{inexecnt}}
% Save the standard \@item before any package (e.g. paralist) redefines it.
\let\@orignalitem\@item

\newlength{\inxlist@labelwidth}
\newlength{\inxlist@labelsep}
\setlength{\inxlist@labelsep}{.75em}
\newlength{\inxlist@itemsep}
\setlength{\inxlist@itemsep}{1.5\p@ plus\p@}

% \theinxcnt: formats the sub-example counter value (letter/roman/etc.)
% Set locally by \inxlist@open to match the current list style.
\def\theinxcnt{\alph{inxcnt}}
%% Shared setup for all inxlist variants.
%% #1 = counter format command, e.g. \alph, \roman, \arabic, \Alph, \Roman
%% #2 = sample string for label width, e.g. "m." or "iii." or "99."
\def\inxlist@open#1#2{%
  \setcounter{inxcnt}{0}%
  \def\inxlist@counterfmt{#1}%
  % Set \theinxcnt so \refstepcounter{inxcnt} records the right label.
  \def\theinxcnt{#1{inxcnt}}%
  \settowidth{\inxlist@labelwidth}{#2}%
  \edef\inxlist@savedlistdepth{\the\@listdepth}%
  \edef\inxlist@saveditemdepth{\the\@itemdepth}%
  \let\@item\@orignalitem
  \begin{list}{}{%
    \labelwidth=\inxlist@labelwidth
    \labelsep=\inxlist@labelsep
    \leftmargin=\inxlist@labelwidth
    \advance\leftmargin\inxlist@labelsep
    \itemindent\z@
    \listparindent\z@
    \topsep=\inxlist@itemsep
    \itemsep=\inxlist@itemsep
    \parsep\z@
  }%
  \def\xi{%
    \refstepcounter{inxcnt}%
    \item[\inxlist@counterfmt{inxcnt}.]%
  }%
  \def\xij[##1]{%
    \refstepcounter{inxcnt}%
    \item[\inxlist@counterfmt{inxcnt}.]##1\hspace{.5ex}%
  }%
  \def\xisn{\item[]}%
  \def\xisnj[##1]{\item[]##1\hspace{.5ex}}%
}

\def\inxlist@close{%
  \end{list}%
  \@listdepth=\inxlist@savedlistdepth\relax
  \@itemdepth=\inxlist@saveditemdepth\relax
  \ignorespacesafterend
}

% inxlist  — alphabetic:  a. b. c.   (mirrors gb4e's xlist / xlista)
\newenvironment{inxlist}
  {\inxlist@open{\alph}{m.}}
  {\inxlist@close}

% inxlisti — roman:       i. ii. iii.  (mirrors gb4e's xlisti)
\newenvironment{inxlisti}
  {\inxlist@open{\roman}{iii.}}
  {\inxlist@close}

% inxlistn — arabic:      1. 2. 3.   (mirrors gb4e's xlistn)
\newenvironment{inxlistn}
  {\inxlist@open{\arabic}{9.}}
  {\inxlist@close}

% inxlistA — Alphabetic:  A. B. C.   (mirrors gb4e's xlistA)
\newenvironment{inxlistA}
  {\inxlist@open{\Alph}{M.}}
  {\inxlist@close}

% inxlistI — Roman:       I. II. III.  (mirrors gb4e's xlistI)
\newenvironment{inxlistI}
  {\inxlist@open{\Roman}{III.}}
  {\inxlist@close}

%% ---------------------------------------------------------------
%% inexe environment
%% ---------------------------------------------------------------

\newif\if@inexe@firstline

\newenvironment{inexe}{%
  \inexe@savedleft=\@totalleftmargin
  \advance\inexe@savedleft\leftskip
  \ifnum\value{inexecnt}<9
    \settowidth{\inexe@labelwidth}{\inexe@fmt{9}}%
  \else\ifnum\value{inexecnt}<99
    \settowidth{\inexe@labelwidth}{\inexe@fmt{99}}%
  \else
    \settowidth{\inexe@labelwidth}{\inexe@fmt{999}}%
  \fi\fi
  \@inexe@firstlinetrue
  \leavevmode\newline\noindent
  \let\ex\inexe@ex
  \let\exj\inexe@exj
  \let\sn\inexe@sn
  \let\snj\inexe@snj
  \let\exi\inexe@exi
}{%
  \vspace{-\inexe@itemsep}%
  \vspace{\belowdisplayskip}%
  \ignorespacesafterend
}

\makeatother

\endinput
